% Generated by tools/build.py -- do not edit; edit data/programming-graphics/*.txt
% No preamble here: this file is \input by main.tex.
\chapter{Unit 3: Programming Languages and Computer Graphics}

\section{Language Design and Translation Issues, Paradigms, Binding and Elementary Data Types}

\begin{question}{H}
Match List I with List II.\\
List I (Programming paradigm): A. Imperative  B. Functional  C. Logic  D. Object-oriented\\
List II (Example language): I. Prolog  II. Haskell  III. C  IV. Smalltalk
\optsshort{A-III, B-II, C-I, D-IV}{A-II, B-III, C-I, D-IV}{A-III, B-II, C-IV, D-I}{A-III, B-I, C-II, D-IV}
\end{question}

\begin{question}{H}
Match List I with List II.\\
List I (Binding time): A. Language design time  B. Language implementation time  C. Compile time  D. Run time\\
List II (Example of what is bound): I. The value of a variable  II. The meaning of the symbol * as multiplication  III. The type of a variable declared int x in C  IV. The range of values of type int (for example 16 or 32 bits)
\optsshort{A-IV, B-II, C-III, D-I}{A-II, B-IV, C-I, D-III}{A-II, B-III, C-IV, D-I}{A-II, B-IV, C-III, D-I}
\end{question}

\begin{question}{H}
Given below are two statements.\\
Statement I: In static typing, type checking is done at compile time.\\
Statement II: In dynamic typing, the type of a variable is bound when a value is assigned at run time.
\optslong{Both Statement I and Statement II are true}{Both Statement I and Statement II are false}{Statement I is true but Statement II is false}{Statement I is false but Statement II is true}
\end{question}

\begin{question}{H}
Given below are two statements: one is labelled as Assertion A and the other is labelled as Reason R.\\
Assertion A: Early (static) binding generally leads to more efficient programs than late (dynamic) binding.\\
Reason R: With early binding, decisions are made once at compile time instead of repeatedly during execution.
\optslong{Both A and R are true and R is the correct explanation of A}{Both A and R are true but R is NOT the correct explanation of A}{A is true but R is false}{A is false but R is true}
\end{question}

\begin{question}{H}
Arrange the stages of translating a high-level program into an executing program in order.\\
A. Loading  B. Compilation  C. Linking  D. Execution
\optsshort{C, B, A, D}{B, A, C, D}{A, B, C, D}{B, C, A, D}
\end{question}

\begin{question}{H}
Which of the following are properties of an imperative language?\\
A. Variables representing memory cells  B. Assignment statements  C. Iteration  D. Referential transparency as the central principle\\
Choose the correct answer from the options given below:
\optsshort{A, B and C only}{A, B, C and D}{B, C and D only}{A and D only}
\end{question}

\begin{question}{M}
A virtual computer, in the sense of Pratt and Zelkowitz, is
\optslong{a physical machine}{the computer as seen by the programmer, defined by the programming language and implemented by hardware, firmware and software}{a virus}{a network computer}
\end{question}

\begin{question}{H}
Match List I with List II.\\
List I (Concept): A. Lifetime  B. Scope  C. l-value  D. r-value\\
List II (Meaning): I. The region of program text in which a name is visible  II. The period of execution during which storage is bound to a variable  III. The value stored in a variable  IV. The address (location) of a variable
\optsshort{A-I, B-II, C-IV, D-III}{A-II, B-I, C-III, D-IV}{A-II, B-IV, C-I, D-III}{A-II, B-I, C-IV, D-III}
\end{question}

\begin{question}{H}
Given below are two statements.\\
Statement I: A strongly typed language detects all type errors.\\
Statement II: C is a strongly typed language because it allows implicit conversions and unions without checks.
\optslong{Both Statement I and Statement II are true}{Both Statement I and Statement II are false}{Statement I is true but Statement II is false}{Statement I is false but Statement II is true}
\end{question}

\begin{question}{M}
Scalar (elementary) data types include
\optslong{integers, reals, characters, Booleans, enumerations and pointers}{lists and sets}{classes and files}{arrays and records}
\end{question}

\begin{question}{M}
Given below are two statements: one is labelled as Assertion A and the other is labelled as Reason R.\\
Assertion A: An interpreter is useful during program development.\\
Reason R: Interpreted programs run faster than compiled programs.
\optslong{Both A and R are true and R is the correct explanation of A}{Both A and R are true but R is NOT the correct explanation of A}{A is true but R is false}{A is false but R is true}
\end{question}

\begin{question}{M}
A language feature that allows a name to be bound to different types or objects at different times is related to
\optsshort{constant folding}{dynamic binding}{static binding}{macro expansion}
\end{question}

\begin{question}{M}
Formal methods for describing programming language syntax include
\optslong{flowcharts only}{BNF and EBNF grammars and syntax diagrams}{UML only}{ER diagrams}
\end{question}

\begin{question}{M}
Formal models of semantics include
\optslong{lexical analysis}{only syntax diagrams}{regular expressions only}{operational, denotational and axiomatic semantics}
\end{question}

\begin{question}{M}
Given below are two statements.\\
Statement I: In C, the type of an expression can be changed with a cast.\\
Statement II: Implicit type conversion (coercion) in C never loses information.
\optslong{Both Statement I and Statement II are true}{Both Statement I and Statement II are false}{Statement I is true but Statement II is false}{Statement I is false but Statement II is true}
\end{question}

\begin{question}{M}
Which of the following is a declarative language?
\optsshort{C}{SQL}{Assembly}{Pascal}
\end{question}

\begin{question}{M}
A programming environment includes
\optslong{only the hardware}{only the compiler}{only the operating system}{tools such as editors, compilers, debuggers, profilers, version control and build tools that support program development}
\end{question}

\begin{question}{M}
Given below are two statements: one is labelled as Assertion A and the other is labelled as Reason R.\\
Assertion A: Java programs are portable across platforms.\\
Reason R: Java provides automatic garbage collection.
\optslong{Both A and R are true and R is the correct explanation of A}{Both A and R are true but R is NOT the correct explanation of A}{A is true but R is false}{A is false but R is true}
\end{question}

\begin{question}{M}
Aliasing occurs when
\optslong{a program uses recursion}{a variable has no name}{two or more names refer to the same memory location}{a name refers to two locations}
\end{question}

\begin{question}{M}
Garbage collection is used to
\optslong{reclaim heap storage that is no longer reachable}{delete files}{clear the stack}{remove dead code}
\end{question}

\begin{question}{M}
Which of the following is a composite (structured) data type?
\optsshort{Character}{Boolean}{Integer}{Record (structure)}
\end{question}

\begin{question}{M}
An enumeration type in C
\optslong{is a user-defined type whose values are named integer constants}{is a floating-point type}{stores strings}{is a pointer type}
\end{question}

\begin{question}{M}
Given below are two statements.\\
Statement I: Static scoping determines non-local references from the program text.\\
Statement II: Most modern languages, including C, Java and Python, use dynamic scoping by default.
\optslong{Both Statement I and Statement II are true}{Both Statement I and Statement II are false}{Statement I is true but Statement II is false}{Statement I is false but Statement II is true}
\end{question}

\begin{question}{L}
FORTRAN, one of the earliest high-level languages, was designed mainly for
\optslong{scientific and numerical computation}{artificial intelligence}{business data processing}{web development}
\end{question}

\begin{question}{M}
Match List I with List II.\\
List I (Language): A. COBOL  B. LISP  C. Prolog  D. Ada\\
List II (Main application area or feature): I. Logic programming and AI  II. Business data processing  III. Safety-critical and embedded systems  IV. Symbolic processing with lists
\optsshort{A-IV, B-II, C-I, D-III}{A-II, B-IV, C-I, D-III}{A-II, B-I, C-IV, D-III}{A-II, B-IV, C-III, D-I}
\end{question}

\begin{question}{M}
A type system is called type-safe if
\optslong{it has no types}{it prevents operations from being applied to values of inappropriate types}{it allows any operation on any value}{it is interpreted}
\end{question}

\begin{question}{M}
Given below are two statements.\\
Statement I: An interpreted program generally executes faster than the same program compiled to machine code.\\
Statement II: Java source code is compiled directly into native machine code for each processor.
\optslong{Both Statement I and Statement II are true}{Both Statement I and Statement II are false}{Statement I is true but Statement II is false}{Statement I is false but Statement II is true}
\end{question}

\begin{question}{M}
Which of the following languages uses indentation to define blocks?
\optsshort{Pascal}{C}{Java}{Python}
\end{question}

\begin{question}{L}
The ability of a language to treat functions as values that can be passed as arguments and returned from functions is called
\optsshort{static binding}{first-class functions}{overloading}{coercion}
\end{question}

\begin{question}{M}
Given below are two statements: one is labelled as Assertion A and the other is labelled as Reason R.\\
Assertion A: Pure functional programs are easier to reason about and parallelise.\\
Reason R: Pure functions have no side effects, so a call always returns the same result for the same arguments.
\optslong{Both A and R are true and R is the correct explanation of A}{Both A and R are true but R is NOT the correct explanation of A}{A is true but R is false}{A is false but R is true}
\end{question}

\begin{answers}
\ans{1}{1}{Imperative languages change state with assignments; functional languages evaluate expressions; logic languages use rules and queries.}
\ans{2}{4}{Early binding gives efficiency; late binding gives flexibility.}
\ans{3}{1}{C, C++ and Java are statically typed; Python and JavaScript are dynamically typed.}
\ans{4}{1}{Late binding, as in dynamic dispatch, gives flexibility at a run-time cost.}
\ans{5}{4}{The compiler produces object code, the linker combines modules and libraries, and the loader places the program in memory.}
\ans{6}{1}{Referential transparency is the hallmark of pure functional languages.}
\ans{7}{2}{For example, the "C virtual computer" executes C programs.}
\ans{8}{4}{In x = x + 1, the left x denotes an l-value and the right x an r-value.}
\ans{9}{3}{C permits unchecked unions, casts and pointer arithmetic, so it is usually considered weakly typed.}
\ans{10}{1}{Composite data types include arrays, records (structures), unions, strings, sets and files.}
\ans{11}{3}{Interpreters are useful because they report errors immediately and allow statement-by-statement execution; compiled programs usually run faster.}
\ans{12}{2}{Dynamic binding of methods to calls is the basis of run-time polymorphism.}
\ans{13}{2}{EBNF adds notation for optional and repeated items.}
\ans{14}{4}{Axiomatic semantics uses pre- and post-conditions (Hoare triples).}
\ans{15}{3}{Conversions such as double to int or long to char can lose information.}
\ans{16}{2}{Declarative languages specify what is wanted, not how to compute it.}
\ans{17}{4}{Integrated development environments (IDEs) combine these tools.}
\ans{18}{2}{Both are true, but portability comes from compiling to bytecode that runs on any JVM.}
\ans{19}{3}{Pointers and reference parameters can create aliases, which make programs harder to reason about.}
\ans{20}{1}{Java, Python and many other languages provide automatic garbage collection.}
\ans{21}{4}{A record groups elements of possibly different types.}
\ans{22}{1}{For example enum colour \{ RED, GREEN, BLUE \}.}
\ans{23}{3}{They use static (lexical) scoping; early Lisp and some shells used dynamic scoping.}
\ans{24}{1}{COBOL was designed for business data processing and LISP for AI.}
\ans{25}{2}{Ada was developed for the US Department of Defense.}
\ans{26}{2}{Type safety rules out errors such as treating an integer as a pointer.}
\ans{27}{2}{Compiled code is generally faster, and Java is compiled to bytecode that is interpreted or JIT-compiled by the JVM.}
\ans{28}{4}{C, Java and many others use braces or keywords such as begin and end.}
\ans{29}{2}{Higher-order functions are central to functional programming.}
\ans{30}{1}{This property is referential transparency.}
\end{answers}

\section{Programming in C: Tokens, Data Types, Operators and Control Statements (with output questions)}

\begin{question}{H}
What is the output of the following C code?\\
int a = 10, b = 3; printf("\%d \%d \%d", a / b, a \% b, -a \% b);
\optsshort{3 -1 -1}{3.33 1 -1}{3 1 -1}{3 1 1}
\end{question}

\begin{question}{H}
What is the output of the following C code?\\
printf("\%d", 5 > 3 > 1);
\optsshort{0}{Compilation error}{5}{1}
\end{question}

\begin{question}{H}
What is the output of the following C code?\\
int i = 0; while (i++ < 3); printf("\%d", i);
\optsshort{3}{0}{Infinite loop}{4}
\end{question}

\begin{question}{H}
What is the output of the following C code?\\
int x = 2; switch (x) \{ case 1: printf("A"); case 2: printf("B"); case 3: printf("C"); break; default: printf("D"); \}
\optsshort{B}{BCD}{BC}{ABC}
\end{question}

\begin{question}{H}
What is the output of the following C code?\\
int i; for (i = 0; i < 5; i++) \{ if (i == 3) continue; printf("\%d", i); \}
\optsshort{01234}{0123}{012}{0124}
\end{question}

\begin{question}{H}
What is the output of the following C code?\\
int x = 2, y = 5; x = x \textasciicircum{} y; y = x \textasciicircum{} y; x = x \textasciicircum{} y; printf("\%d \%d", x, y);
\optsshort{5 2}{7 7}{2 5}{0 0}
\end{question}

\begin{question}{H}
What is the output of the following C code?\\
printf("\%d \%d \%d \%d", 5 \& 3, 5 | 3, 5 \textasciicircum{} 3, 1 << 4);
\optsshort{1 6 7 16}{1 7 6 16}{7 1 6 16}{1 7 6 8}
\end{question}

\begin{question}{H}
What is the output of the following C code?\\
int i = 0, j = 0; if (i++ || j++) ; printf("\%d \%d", i, j);
\optsshort{1 1}{0 0}{0 1}{1 0}
\end{question}

\begin{question}{H}
What is the output of the following C code?\\
int i = 0, j = 0; if (++i || ++j) ; printf("\%d \%d", i, j);
\optsshort{1 1}{0 0}{1 0}{0 1}
\end{question}

\begin{question}{M}
What is the output of the following C code?\\
int i = 0, j = 0; if (i++ \&\& j++) ; printf("\%d \%d", i, j);
\optsshort{0 0}{1 0}{1 1}{0 1}
\end{question}

\begin{question}{H}
What is the output of the following C code?\\
int x = -1; unsigned int y = 1; printf("\%d", x < y);
\optsshort{1}{-1}{Compilation error}{0}
\end{question}

\begin{question}{H}
What is the output of the following C code?\\
float f = 0.1; if (f == 0.1) printf("Yes"); else printf("No");
\optsshort{No}{Yes}{Compilation error}{0.1}
\end{question}

\begin{question}{H}
What is the output of the following C code?\\
enum \{ A, B = 5, C, D = 2, E \}; printf("\%d \%d \%d", A, C, E);
\optsshort{0 2 3}{0 6 2}{0 6 3}{1 6 3}
\end{question}

\begin{question}{M}
What is the output of the following C code?\\
int i = 10; do \{ printf("\%d", i); i++; \} while (i < 5);
\optsshort{10}{Nothing}{Infinite loop}{10111213}
\end{question}

\begin{question}{H}
What is the output of the following C code?\\
printf("\%d \%d", -7 / 2, -7 \% 2);
\optsshort{-3 1}{-3 -1}{-4 -1}{-4 1}
\end{question}

\begin{question}{M}
What is the output of the following C code?\\
int a = (1, 2, 3); printf("\%d", a);
\optsshort{Compilation error}{6}{3}{1}
\end{question}

\begin{question}{H}
What is the output of the following C code?\\
printf("\%d", printf("Hello"));
\optsshort{Hello6}{Hello}{Hello5}{5Hello}
\end{question}

\begin{question}{M}
In C, the value of sizeof('A') on a system with 4-byte int is
\optsshort{8}{2}{1}{4}
\end{question}

\begin{question}{M}
What is the output of the following C code?\\
int i, c = 0; for (i = 1; i <= 10; i *= 2) c++; printf("\%d", c);
\optsshort{4}{10}{5}{3}
\end{question}

\begin{question}{M}
What is the output of the following C code?\\
int i = 5, j; j = i++ * 2; printf("\%d \%d", i, j);
\optsshort{5 10}{6 12}{6 10}{5 12}
\end{question}

\begin{question}{M}
What is the output of the following C code?\\
int a = 5, b = 10, c = 3; printf("\%d", a > b ? a : b > c ? b : c);
\optsshort{1}{3}{5}{10}
\end{question}

\begin{question}{M}
What is the output of the following C code?\\
int x = 2 + 3 * 4 \% 5 - 6 / 4; printf("\%d", x);
\optsshort{4}{3}{1}{5}
\end{question}

\begin{question}{M}
What is the output of the following C code?\\
int x = 0; if (x = 5) printf("T\%d", x); else printf("F\%d", x);
\optsshort{F0}{T5}{T0}{F5}
\end{question}

\begin{question}{M}
What is the output of the following C code?\\
int n = 0; while (1) \{ n += 3; if (n > 10) break; \} printf("\%d", n);
\optsshort{10}{9}{Infinite loop}{12}
\end{question}

\begin{question}{M}
What is the output of the following C code?\\
int i, j, c = 0; for (i = 0; i < 4; i++) for (j = i; j < 4; j++) c++; printf("\%d", c);
\optsshort{12}{8}{10}{16}
\end{question}

\begin{question}{M}
What is the output of the following C code?\\
char c = 'A' + 2; printf("\%c \%d", c, c);
\optsshort{C 67}{B 66}{C 65}{A 65}
\end{question}

\begin{question}{H}
Given below are two statements.\\
Statement I: In C, the expression i = i++ + ++i has undefined behaviour.\\
Statement II: In C, the order of evaluation of function arguments is unspecified.
\optslong{Both Statement I and Statement II are true}{Both Statement I and Statement II are false}{Statement I is true but Statement II is false}{Statement I is false but Statement II is true}
\end{question}

\begin{question}{H}
Match List I with List II.\\
List I (C operator): A. ++ (postfix)  B. * (multiplication)  C. \&\& (logical AND)  D. = (assignment)\\
List II (Associativity / relative precedence): I. Lower precedence than relational operators, left to right  II. Highest of these, postfix  III. Lowest of these, right to left  IV. Higher than + and $-$, left to right
\optsshort{A-II, B-IV, C-III, D-I}{A-IV, B-II, C-I, D-III}{A-II, B-IV, C-I, D-III}{A-II, B-I, C-IV, D-III}
\end{question}

\begin{question}{M}
Which of the following is NOT a valid C identifier?
\optsshort{\_\allowbreak{}count}{count2}{Count\_\allowbreak{}2}{2count}
\end{question}

\begin{question}{M}
Which of the following are C keywords?\\
A. auto  B. register  C. volatile  D. include  E. typedef\\
Choose the correct answer from the options given below:
\optsshort{A, B, C, D and E}{C, D and E only}{A, B and D only}{A, B, C and E only}
\end{question}

\begin{question}{M}
Given below are two statements: one is labelled as Assertion A and the other is labelled as Reason R.\\
Assertion A: The break statement exits only the innermost loop or switch that contains it.\\
Reason R: C has no labelled break; goto must be used to exit several nested loops at once.
\optslong{Both A and R are true and R is the correct explanation of A}{Both A and R are true but R is NOT the correct explanation of A}{A is true but R is false}{A is false but R is true}
\end{question}

\begin{question}{M}
The goto statement in C transfers control to
\optslong{any function}{a labelled statement within the same function}{the beginning of the program}{the operating system}
\end{question}

\begin{question}{M}
Given below are two statements.\\
Statement I: In C, a switch expression can be of type float.\\
Statement II: Case labels in a C switch statement must be integer constant expressions.
\optslong{Both Statement I and Statement II are true}{Both Statement I and Statement II are false}{Statement I is true but Statement II is false}{Statement I is false but Statement II is true}
\end{question}

\begin{question}{M}
The range of values of a signed 8-bit char in two's complement representation is
\optsshort{$-$127 to 127}{$-$256 to 255}{$-$128 to 127}{0 to 255}
\end{question}

\begin{question}{L}
Which format specifier prints an unsigned integer in hexadecimal in C?
\optsshort{\%o}{\%d}{\%f}{\%x}
\end{question}

\begin{question}{M}
Given below are two statements: one is labelled as Assertion A and the other is labelled as Reason R.\\
Assertion A: The expression 5 / 2 in C evaluates to 2.5.\\
Reason R: When both operands are integers, C performs integer division.
\optslong{Both A and R are true and R is the correct explanation of A}{Both A and R are true but R is NOT the correct explanation of A}{A is true but R is false}{A is false but R is true}
\end{question}

\begin{answers}
\ans{1}{3}{Integer division truncates towards zero, and the remainder takes the sign of the dividend (C99).}
\ans{2}{1}{Relational operators associate left to right: (5 > 3) is 1, and 1 > 1 is 0.}
\ans{3}{4}{The loop body is the empty statement; the test is evaluated with i = 0, 1, 2, 3, and i is incremented each time, ending at 4.}
\ans{4}{3}{Control falls through from case 2 into case 3 until the break.}
\ans{5}{4}{continue skips the printf only when i is 3.}
\ans{6}{1}{The three XOR assignments swap x and y without a temporary variable.}
\ans{7}{2}{101 \& 011 = 001, 101 | 011 = 111, 101 \textasciicircum{} 011 = 110, and 1 << 4 = 16.}
\ans{8}{1}{i++ yields 0 (false), so the right operand j++ is also evaluated.}
\ans{9}{3}{++i yields 1 (true), so ++j is not evaluated (short-circuit evaluation).}
\ans{10}{2}{i++ yields 0, so the \&\& short-circuits and j++ is skipped.}
\ans{11}{4}{x is converted to unsigned int, becoming a very large value, so x < y is false.}
\ans{12}{1}{0.1 cannot be represented exactly; the float value differs from the double constant 0.1.}
\ans{13}{3}{Enumerators take the previous value plus one unless explicitly assigned.}
\ans{14}{1}{A do-while loop executes its body at least once.}
\ans{15}{2}{C99 truncates towards zero, so a = (a / b) $\times$ b + a \% b gives -7 = -3 $\times$ 2 + (-1).}
\ans{16}{3}{The comma operator evaluates its operands left to right and yields the value of the last one.}
\ans{17}{3}{The inner printf prints Hello and returns the number of characters printed, 5.}
\ans{18}{4}{A character constant has type int in C (in C++ it has type char, so the size is 1).}
\ans{19}{1}{i takes the values 1, 2, 4 and 8.}
\ans{20}{3}{Post-increment uses the old value 5 in the expression and then increments i.}
\ans{21}{4}{?: is right-associative: a > b is false, so b > c ? b : c gives 10.}
\ans{22}{2}{* , \% and / have equal precedence and are left-associative: 3 $\times$ 4 = 12, 12 \% 5 = 2, 6 / 4 = 1; 2 + 2 $-$ 1 = 3.}
\ans{23}{2}{x = 5 is an assignment, whose value 5 is true.}
\ans{24}{4}{n takes the values 3, 6, 9 and 12, and the loop exits when n exceeds 10.}
\ans{25}{3}{4 + 3 + 2 + 1 = 10.}
\ans{26}{1}{Characters are small integers; 'A' is 65 in ASCII.}
\ans{27}{1}{Modifying a variable twice without an intervening sequence point is undefined behaviour.}
\ans{28}{3}{Assignment operators have very low precedence and associate right to left.}
\ans{29}{4}{Identifiers cannot begin with a digit.}
\ans{30}{4}{include is a preprocessor directive name, not a keyword.}
\ans{31}{1}{Java, in contrast, provides labelled break statements.}
\ans{32}{2}{Structured programming discourages goto, but it is occasionally used for error handling.}
\ans{33}{4}{The switch expression must have integer type (including char and enum).}
\ans{34}{3}{An n-bit two's complement integer ranges from $-$2\textsuperscript{n$-$1} to 2\textsuperscript{n$-$1} $-$ 1.}
\ans{35}{4}{\%o prints octal and \%u unsigned decimal.}
\ans{36}{4}{Because of R, 5 / 2 evaluates to 2.}
\end{answers}

\section{Programming in C: Arrays, Strings, Pointers, Functions, Structures, Unions, Storage Classes, Preprocessor, Files and Command-Line Arguments}

\begin{question}{H}
What is the output of the following C code?\\
char s[] = "hello"; printf("\%zu \%zu", sizeof(s), strlen(s));
\optsshort{6 6}{6 5}{5 6}{5 5}
\end{question}

\begin{question}{H}
What is the output of the following C code?\\
int a[] = \{1, 2, 3, 4, 5\}; int *p = a; p += 2; printf("\%d \%d", *p, *(p + 2));
\optsshort{1 3}{3 4}{2 4}{3 5}
\end{question}

\begin{question}{H}
On a system with 4-byte int, what is the output of the following C code?\\
int a[3][4]; printf("\%zu", sizeof(a));
\optsshort{48}{16}{8}{12}
\end{question}

\begin{question}{H}
What is the output of the following C code?\\
char *p = "ABCDE"; printf("\%c \%s", *(p + 2), p + 3);
\optsshort{C CDE}{C DE}{B DE}{D DE}
\end{question}

\begin{question}{H}
What is the output of the following C code?\\
int x = 10; int *p = \&x; int **q = \&p; **q = 20; printf("\%d", x);
\optsshort{20}{Address of x}{Compilation error}{10}
\end{question}

\begin{question}{H}
What is the output of the following C code?\\
int a[] = \{10, 20, 30\}; printf("\%d \%d", *a + 1, *(a + 1));
\optsshort{20 11}{11 11}{20 20}{11 20}
\end{question}

\begin{question}{H}
What is the output of the following C code?\\
int a[2][3] = \{\{1, 2, 3\}, \{4, 5, 6\}\}; int (*p)[3] = a; printf("\%d \%d", *(*(p + 1) + 2), **p);
\optsshort{5 1}{6 4}{3 1}{6 1}
\end{question}

\begin{question}{H}
What is the output of the following C code?\\
\#define SQ(x) x * x\\
printf("\%d", SQ(3 + 2));
\optsshort{11}{25}{13}{10}
\end{question}

\begin{question}{H}
What is the output of the following C code, in which f is defined as int f(void) \{ static int c = 0; return ++c; \}?\\
f(); f(); printf("\%d", f());
\optsshort{0}{3}{1}{2}
\end{question}

\begin{question}{H}
What is the output of the following C code?\\
void f(int n) \{ if (n > 0) \{ f(n - 1); printf("\%d", n); \} \}  --- called as f(3)
\optsshort{3210}{321}{123}{0123}
\end{question}

\begin{question}{H}
What is the output of f(1234), where int f(int n) \{ if (n == 0) return 0; return n \% 10 + f(n / 10); \}?
\optsshort{4}{24}{10}{1234}
\end{question}

\begin{question}{H}
What is the output of the following C code?\\
void swap(int a, int b) \{ int t = a; a = b; b = t; \}  --- main: int x = 1, y = 2; swap(x, y); printf("\%d \%d", x, y);
\optsshort{1 1}{1 2}{2 2}{2 1}
\end{question}

\begin{question}{H}
On a typical system where int is 4 bytes with 4-byte alignment, sizeof(struct s) for struct s \{ char c; int i; char d; \} is
\optsshort{8}{6}{12}{16}
\end{question}

\begin{question}{M}
On a typical system where int is 4 bytes with 4-byte alignment, sizeof(struct s) for struct s \{ int i; char c; char d; \} is
\optsshort{12}{8}{6}{10}
\end{question}

\begin{question}{H}
On a typical system with 4-byte int and 8-byte double, sizeof(union u) for union u \{ int i; char c[6]; double d; \} is
\optsshort{6}{18}{8}{24}
\end{question}

\begin{question}{H}
Given below are two statements.\\
Statement I: All members of a union share the same memory location.\\
Statement II: All members of a structure can hold valid values at the same time.
\optslong{Both Statement I and Statement II are true}{Both Statement I and Statement II are false}{Statement I is true but Statement II is false}{Statement I is false but Statement II is true}
\end{question}

\begin{question}{H}
On a 64-bit system with 4-byte int, the code void g(int a[]) \{ printf("\%zu ", sizeof(a)); \} with main: int a[10]; g(a); printf("\%zu", sizeof(a)); prints
\optsshort{40 40}{40 8}{4 40}{8 40}
\end{question}

\begin{question}{H}
Match List I with List II.\\
List I (Declaration): A. int *p[10];  B. int (*p)[10];  C. int (*p)(int);  D. int *p(int);\\
List II (Meaning): I. A function taking int and returning a pointer to int  II. An array of 10 pointers to int  III. A pointer to an array of 10 ints  IV. A pointer to a function taking int and returning int
\optsshort{A-II, B-IV, C-III, D-I}{A-III, B-II, C-IV, D-I}{A-II, B-III, C-I, D-IV}{A-II, B-III, C-IV, D-I}
\end{question}

\begin{question}{H}
Match List I with List II.\\
List I (Storage class): A. auto  B. register  C. static (local)  D. extern\\
List II (Property): I. Declares a variable defined in another file or later in the program  II. Default for local variables; lifetime of the block  III. Suggests keeping the variable in a CPU register; its address cannot be taken  IV. Lifetime of the whole program but scope of the block
\optsshort{A-II, B-III, C-I, D-IV}{A-II, B-IV, C-III, D-I}{A-III, B-II, C-IV, D-I}{A-II, B-III, C-IV, D-I}
\end{question}

\begin{question}{M}
What is the output of the following C code?\\
char s1[] = "abc", s2[] = "abc"; printf("\%d \%d", s1 == s2, strcmp(s1, s2));
\optsshort{0 1}{1 1}{0 0}{1 0}
\end{question}

\begin{question}{M}
What is the output of the following C code?\\
printf("\%zu", strlen("abc\textbackslash{}0def"));
\optsshort{8}{3}{6}{7}
\end{question}

\begin{question}{M}
What is the output of the following C code?\\
\#define STR(x) \#x\\
\#define CAT(a, b) a\#\#b\\
int xy = 7; printf("\%s \%d", STR(hello), CAT(x, y));
\optsshort{Compilation error}{x y 7}{hello 7}{"hello" xy}
\end{question}

\begin{question}{M}
Given below are two statements: one is labelled as Assertion A and the other is labelled as Reason R.\\
Assertion A: A function-like macro can produce unexpected results when its argument has side effects, such as MAX(i++, j).\\
Reason R: Macro names are conventionally written in capital letters.
\optslong{Both A and R are true and R is the correct explanation of A}{Both A and R are true but R is NOT the correct explanation of A}{A is true but R is false}{A is false but R is true}
\end{question}

\begin{question}{H}
Match List I with List II.\\
List I (fopen mode): A. "r"  B. "w"  C. "a"  D. "r+"\\
List II (Behaviour): I. Opens an existing file for reading and writing  II. Opens for reading; fails if the file does not exist  III. Creates a file or truncates an existing one, for writing  IV. Opens for writing at the end of the file, creating it if necessary
\optsshort{A-II, B-III, C-IV, D-I}{A-III, B-II, C-IV, D-I}{A-II, B-III, C-I, D-IV}{A-II, B-IV, C-III, D-I}
\end{question}

\begin{question}{M}
To find the size of a file in C, a program can use
\optslong{fgetc once}{fopen only}{sizeof(fp)}{fseek(fp, 0, SEEK\_\allowbreak{}END) followed by ftell(fp)}
\end{question}

\begin{question}{H}
A C program is run as ./prog one two three. The values of argc and argv[1] are
\optsshort{3 and "prog"}{3 and "one"}{4 and "./prog"}{4 and "one"}
\end{question}

\begin{question}{M}
Given below are two statements.\\
Statement I: The name of an array in C is a modifiable pointer variable.\\
Statement II: In most expressions, the name of an array decays to a pointer to its first element.
\optslong{Both Statement I and Statement II are true}{Both Statement I and Statement II are false}{Statement I is true but Statement II is false}{Statement I is false but Statement II is true}
\end{question}

\begin{question}{M}
Dereferencing a pointer that has been freed is an example of
\optslong{a dangling pointer}{a wild pointer that was never initialised}{a null pointer}{a memory leak}
\end{question}

\begin{question}{M}
The function calloc(n, size) differs from malloc(n $\times$ size) in that calloc
\optslong{frees memory automatically}{initialises the allocated memory to zero}{allocates on the stack}{cannot fail}
\end{question}

\begin{question}{M}
Given below are two statements.\\
Statement I: A void pointer can be dereferenced directly in C.\\
Statement II: A void pointer can hold the address of an object of any data type.
\optslong{Both Statement I and Statement II are true}{Both Statement I and Statement II are false}{Statement I is true but Statement II is false}{Statement I is false but Statement II is true}
\end{question}

\begin{question}{M}
In C, a function that calls itself is called
\optsshort{a callback}{a static function}{a recursive function}{an inline function}
\end{question}

\begin{question}{M}
Which header file declares the functions malloc and free?
\optsshort{stdlib.h}{string.h}{stdio.h}{malloc.c}
\end{question}

\begin{question}{M}
Given below are two statements: one is labelled as Assertion A and the other is labelled as Reason R.\\
Assertion A: Passing a large structure by pointer is more efficient than passing it by value.\\
Reason R: Passing by value copies the entire structure.
\optslong{Both A and R are true and R is the correct explanation of A}{Both A and R are true but R is NOT the correct explanation of A}{A is true but R is false}{A is false but R is true}
\end{question}

\begin{question}{M}
The \#ifndef / \#define / \#endif idiom at the top of a header file is used to
\optslong{declare global variables}{speed up linking}{prevent the header from being included more than once in a translation unit (include guard)}{define macros}
\end{question}

\begin{question}{L}
Which function reads a whole line of text safely from a file stream in C?
\optsshort{getch}{gets}{fgets}{scanf("\%s")}
\end{question}

\begin{question}{M}
What does the expression *p++ do when p points into an int array?
\optslong{Increments the value pointed to}{Advances p and then returns the new value}{Is a syntax error}{Returns the value pointed to and then advances p to the next element}
\end{question}

\begin{answers}
\ans{1}{2}{sizeof includes the terminating null character; strlen does not.}
\ans{2}{4}{After p += 2, p points to a[2].}
\ans{3}{1}{3 $\times$ 4 elements $\times$ 4 bytes = 48.}
\ans{4}{2}{*(p + 2) is the third character; p + 3 points to the substring starting at D.}
\ans{5}{1}{**q dereferences twice and refers to x itself.}
\ans{6}{4}{*a + 1 adds 1 to a[0]; *(a + 1) is a[1].}
\ans{7}{4}{p + 1 points to the second row; adding 2 selects its third element, 6.}
\ans{8}{1}{The macro expands to 3 + 2 * 3 + 2 = 11; it should be written as ((x) * (x)).}
\ans{9}{2}{A static local variable keeps its value between calls.}
\ans{10}{3}{Each call prints after its recursive call returns, so the smallest value is printed first.}
\ans{11}{3}{It sums the decimal digits: 4 + 3 + 2 + 1 = 10.}
\ans{12}{2}{C passes arguments by value; swapping requires passing pointers.}
\ans{13}{3}{Padding: c (1) + 3 bytes, i (4), d (1) + 3 bytes of trailing padding.}
\ans{14}{2}{Reordering members so that the small ones are together reduces padding: 4 + 1 + 1 + 2.}
\ans{15}{3}{A union is as large as its largest member (rounded up for alignment).}
\ans{16}{1}{Writing one union member overwrites the others.}
\ans{17}{4}{An array parameter is adjusted to a pointer, so sizeof gives the pointer size (8).}
\ans{18}{4}{[] and () bind more tightly than *, so parentheses change the meaning.}
\ans{19}{4}{C storage classes.}
\ans{20}{3}{s1 == s2 compares the addresses of two different arrays; strcmp compares contents and returns 0 for equal strings.}
\ans{21}{2}{strlen stops at the first null character.}
\ans{22}{3}{\# turns the argument into a string literal; \#\# pastes x and y into the identifier xy.}
\ans{23}{2}{Both are true, but the problem arises because the preprocessor substitutes the argument text, which may then be evaluated more than once.}
\ans{24}{1}{"w+" and "a+" also allow reading.}
\ans{25}{4}{ftell returns the current file position.}
\ans{26}{4}{argv[0] is the program name, so argc counts it too.}
\ans{27}{4}{An array name is not a modifiable lvalue; a = p is an error.}
\ans{28}{1}{A memory leak is allocated memory that is never freed.}
\ans{29}{2}{realloc changes the size of a previously allocated block.}
\ans{30}{4}{A void pointer must be cast to a specific pointer type before it is dereferenced.}
\ans{31}{3}{Each call gets its own activation record on the stack.}
\ans{32}{1}{string.h declares strlen, strcpy and strcmp.}
\ans{33}{1}{A pointer to const can be used when the function must not modify it.}
\ans{34}{3}{\#pragma once is a common non-standard alternative.}
\ans{35}{3}{gets was removed from C11 because it cannot limit the input length.}
\ans{36}{4}{Postfix ++ binds more tightly than *, so it applies to p.}
\end{answers}

\section{Object-Oriented Programming and C++ (classes, constructors, overloading, inheritance, virtual functions, templates, exceptions, streams)}

\begin{question}{H}
What is the output of the following C++ program?\\
class A \{ public: A() \{ cout << "A"; \} \textasciitilde{}A() \{ cout << "a"; \} \};\\
class B : public A \{ public: B() \{ cout << "B"; \} \textasciitilde{}B() \{ cout << "b"; \} \};\\
int main() \{ B obj; \}
\optsshort{ABab}{ABba}{BAab}{BAba}
\end{question}

\begin{question}{H}
What is the output of the following C++ program, where M prints "M"/"m" in its constructor/destructor?\\
class B : public A \{ M m; public: B() \{ cout << "B"; \} \textasciitilde{}B() \{ cout << "b"; \} \};  --- A prints "A"/"a"; main: B obj;
\optsshort{AMBabm}{ABMmba}{AMBbma}{MABbam}
\end{question}

\begin{question}{H}
What is the output of the following C++ code?\\
class Base \{ public: virtual void show() \{ cout << "Base"; \} \};\\
class Derived : public Base \{ public: void show() \{ cout << "Derived"; \} \};\\
Base *p = new Derived; p->show();
\optsshort{Compilation error}{BaseDerived}{Derived}{Base}
\end{question}

\begin{question}{H}
In the previous program, if show() in Base is NOT declared virtual, the call p->show() through a Base pointer prints
\optsshort{Derived}{Compilation error}{BaseDerived}{Base}
\end{question}

\begin{question}{H}
What is the output of the following C++ code?\\
class B \{ public: \textasciitilde{}B() \{ cout << "\textasciitilde{}B"; \} \};\\
class D : public B \{ public: \textasciitilde{}D() \{ cout << "\textasciitilde{}D"; \} \};\\
B *p = new D; delete p;
\optsshort{\textasciitilde{}D}{\textasciitilde{}D\textasciitilde{}B}{\textasciitilde{}B}{\textasciitilde{}B\textasciitilde{}D}
\end{question}

\begin{question}{H}
Given below are two statements: one is labelled as Assertion A and the other is labelled as Reason R.\\
Assertion A: A base class intended for polymorphic use should have a virtual destructor.\\
Reason R: Deleting a derived object through a base pointer without a virtual destructor fails to call the derived destructor.
\optslong{Both A and R are true and R is the correct explanation of A}{Both A and R are true but R is NOT the correct explanation of A}{A is true but R is false}{A is false but R is true}
\end{question}

\begin{question}{H}
In C++, the value printed by cout << sizeof('A'); and by cout << sizeof(E); for an empty class E \{\}; are respectively
\optsshort{4 and 1}{1 and 1}{4 and 0}{1 and 0}
\end{question}

\begin{question}{H}
What is the output of the following C++ code?\\
void f(int x) \{ cout << "int "; \}  void f(double x) \{ cout << "double "; \}\\
int main() \{ f(5); f(5.0); f('a'); \}
\optsshort{int int int}{int double int}{int double double}{Compilation error}
\end{question}

\begin{question}{H}
What is the output of the following C++ code?\\
int add(int a, int b = 10, int c = 20) \{ return a + b + c; \}\\
cout << add(1) << " " << add(1, 2) << " " << add(1, 2, 3);
\optsshort{31 23 33}{31 23 6}{1 3 6}{31 13 6}
\end{question}

\begin{question}{H}
What is the output of the following C++ code?\\
class C \{ public: static int n; C() \{ n++; \} \};  int C::n = 0;\\
int main() \{ C a, b, c; cout << C::n; \}
\optsshort{1}{3}{0}{Compilation error}
\end{question}

\begin{question}{H}
What is the output of the following C++ code?\\
try \{ throw 5; cout << "X"; \} catch (int e) \{ cout << "int " << e; \} catch (...) \{ cout << "other"; \}
\optsshort{X int 5}{int 5}{int 5 other}{other}
\end{question}

\begin{question}{M}
What is the output of the following C++ code?\\
try \{ throw 'a'; \} catch (int e) \{ cout << "int"; \} catch (...) \{ cout << "other"; \}
\optsshort{97}{other}{a}{int}
\end{question}

\begin{question}{H}
What is the output of the following C++ code?\\
class C \{ public: C() \{\} C(const C \&) \{ cout << "copy "; \} \};\\
void f(C c) \{\}\\
int main() \{ C a; f(a); C b = a; \}
\optsshort{copy}{nothing}{copy copy}{copy copy copy}
\end{question}

\begin{question}{M}
What is the output of the following C++ code?\\
int x = 1;\\
int main() \{ int x = 2; cout << x << ::x; \}
\optsshort{12}{21}{22}{11}
\end{question}

\begin{question}{M}
What is the output of the following C++ code?\\
int a = 5; int \&r = a; r = 10; cout << a;
\optsshort{5}{Address of a}{Compilation error}{10}
\end{question}

\begin{question}{M}
What is the output of the following C++ code?\\
template <class T> T mx(T a, T b) \{ return a > b ? a : b; \}\\
cout << mx(3, 7) << " " << mx(2.5, 1.5);
\optsshort{7 2}{7 2.5}{3 1.5}{Compilation error}
\end{question}

\begin{question}{M}
What is the output of the following C++ code?\\
class P \{ public: int x; P(int v) : x(v) \{\} P operator+(const P \&o) \{ return P(x + o.x * 2); \} \};\\
P a(3), b(4); P c = a + b; cout << c.x;
\optsshort{7}{11}{14}{10}
\end{question}

\begin{question}{M}
In C++, a class has prefix operator++ adding 10 and postfix operator++(int) adding 1 to a member v, initially 0. After ++a; a++; the value of a.v is
\optsshort{20}{11}{2}{10}
\end{question}

\begin{question}{H}
Which of the following operators CANNOT be overloaded in C++?\\
A. ::  B. .  C. ?:  D. sizeof  E. []\\
Choose the correct answer from the options given below:
\optsshort{A, B, C, D and E}{A, B, C and D only}{A and B only}{C, D and E only}
\end{question}

\begin{question}{H}
Match List I with List II.\\
List I (OOP concept): A. Encapsulation  B. Inheritance  C. Polymorphism  D. Abstraction\\
List II (Meaning): I. One interface with many implementations, chosen by the object's type  II. Bundling data with the methods that operate on it and hiding internal details  III. Showing only essential features and hiding implementation complexity  IV. A class acquiring the properties of another class
\optsshort{A-IV, B-II, C-I, D-III}{A-II, B-I, C-IV, D-III}{A-II, B-IV, C-III, D-I}{A-II, B-IV, C-I, D-III}
\end{question}

\begin{question}{H}
Given below are two statements.\\
Statement I: A class that has at least one pure virtual function is an abstract class and cannot be instantiated.\\
Statement II: A derived class of an abstract class is abstract unless it overrides all the pure virtual functions.
\optslong{Both Statement I and Statement II are true}{Both Statement I and Statement II are false}{Statement I is true but Statement II is false}{Statement I is false but Statement II is true}
\end{question}

\begin{question}{H}
Match List I with List II.\\
List I (Base-class member under public inheritance): A. public  B. protected  C. private\\
List II (Access in the derived class): I. Not accessible  II. Remains protected  III. Remains public
\optsshort{A-III, B-I, C-II}{A-III, B-II, C-I}{A-I, B-II, C-III}{A-II, B-III, C-I}
\end{question}

\begin{question}{M}
The diamond problem in multiple inheritance (D derives from B and C, which both derive from A) is resolved in C++ by
\optslong{using templates}{declaring B and C with virtual inheritance from A}{using friend functions}{making A abstract}
\end{question}

\begin{question}{M}
Given below are two statements: one is labelled as Assertion A and the other is labelled as Reason R.\\
Assertion A: A friend function of a class can access its private members.\\
Reason R: A friend function is a member function of that class.
\optslong{Both A and R are true and R is the correct explanation of A}{Both A and R are true but R is NOT the correct explanation of A}{A is true but R is false}{A is false but R is true}
\end{question}

\begin{question}{M}
Function overloading in C++ is an example of
\optslong{compile-time (static) polymorphism}{run-time polymorphism}{dynamic binding}{encapsulation}
\end{question}

\begin{question}{M}
The this pointer in a non-static member function points to
\optslong{the first data member only}{the base class}{the class}{the object on which the function was invoked}
\end{question}

\begin{question}{M}
Given below are two statements.\\
Statement I: A constructor can be declared virtual in C++.\\
Statement II: A destructor can be declared virtual in C++.
\optslong{Both Statement I and Statement II are true}{Both Statement I and Statement II are false}{Statement I is true but Statement II is false}{Statement I is false but Statement II is true}
\end{question}

\begin{question}{M}
Which of the following statements about new and malloc in C++ is correct?
\optslong{new returns void* that must be cast}{delete can free memory obtained from malloc}{malloc calls constructors}{new allocates memory and calls the constructor; delete calls the destructor and frees memory}
\end{question}

\begin{question}{M}
Match List I with List II.\\
List I (C++ stream object or class): A. cin  B. cerr  C. ofstream  D. ifstream\\
List II (Purpose): I. Unbuffered standard error output  II. Standard input  III. Reading from a file  IV. Writing to a file
\optsshort{A-I, B-II, C-IV, D-III}{A-II, B-I, C-IV, D-III}{A-II, B-I, C-III, D-IV}{A-II, B-IV, C-I, D-III}
\end{question}

\begin{question}{M}
An inline function in C++
\optslong{cannot have parameters}{is a request to the compiler to substitute the function body at the call site, avoiding call overhead}{is always expanded at the call site by the compiler}{must be recursive}
\end{question}

\begin{question}{M}
Given below are two statements.\\
Statement I: A template function is type-checked when it is instantiated with specific types.\\
Statement II: Templates provide run-time polymorphism.
\optslong{Both Statement I and Statement II are true}{Both Statement I and Statement II are false}{Statement I is true but Statement II is false}{Statement I is false but Statement II is true}
\end{question}

\begin{question}{M}
In C++, a program can be split into several source files compiled separately. A global variable defined in one file is used in another by
\optslong{using a friend declaration}{defining it again}{using static}{redeclaring it with extern in the other file}
\end{question}

\begin{question}{L}
Which keyword is used to handle an exception in C++?
\optsshort{try}{catch}{finally}{throw}
\end{question}

\begin{question}{M}
Given below are two statements: one is labelled as Assertion A and the other is labelled as Reason R.\\
Assertion A: In C++, a member function declared const cannot modify the object's non-mutable data members.\\
Reason R: A const member function can be called only on non-const objects.
\optslong{Both A and R are true and R is the correct explanation of A}{Both A and R are true but R is NOT the correct explanation of A}{A is true but R is false}{A is false but R is true}
\end{question}

\begin{question}{M}
An exception thrown in a function that has no matching handler anywhere in the call chain results in
\optsshort{the exception being ignored}{a compile-time error}{a call to std::terminate}{the program continuing}
\end{question}

\begin{question}{M}
Event handling in C++ GUI frameworks (such as Qt signals and slots, or callbacks in other libraries) typically relies on
\optslong{the preprocessor only}{function pointers, functors or virtual functions invoked when an event occurs}{exception handling}{file streams}
\end{question}

\begin{question}{M}
What does the following C++ code print? class S \{ public: virtual double area() = 0; \}; class Sq : public S \{ double s; public: Sq(double x) : s(x) \{\} double area() \{ return s * s; \} \};  S *p = new Sq(3); cout << p->area();
\optsshort{6}{Compilation error}{9}{3}
\end{question}

\begin{answers}
\ans{1}{2}{Base constructors run before derived ones; destructors run in the reverse order.}
\ans{2}{3}{Order of construction: base class, then members, then the body of the derived constructor; destruction is the reverse.}
\ans{3}{3}{A virtual function is bound at run time to the type of the object pointed to.}
\ans{4}{4}{Without virtual, the call is bound at compile time using the pointer's static type.}
\ans{5}{3}{Without a virtual destructor, only the base destructor runs (formally the behaviour is undefined); with a virtual destructor the output is \textasciitilde{}D\textasciitilde{}B.}
\ans{6}{1}{Resources held by the derived part would otherwise leak.}
\ans{7}{2}{A char literal has type char in C++; an empty class has size 1 so that distinct objects have distinct addresses.}
\ans{8}{2}{'a' is promoted from char to int, which is a better match than conversion to double.}
\ans{9}{2}{Missing trailing arguments take their default values.}
\ans{10}{2}{A static data member is shared by all objects of the class.}
\ans{11}{2}{Control leaves the try block at the throw, and the first matching handler runs.}
\ans{12}{2}{No conversion from char to int is applied when matching handlers, so the catch-all handler runs.}
\ans{13}{3}{The copy constructor is called when an object is passed by value and when an object is initialised from another.}
\ans{14}{2}{::x refers to the global x, hidden by the local x.}
\ans{15}{4}{A reference is an alias for the variable it is bound to.}
\ans{16}{2}{The template is instantiated separately for int and for double.}
\ans{17}{2}{a + b calls a.operator+(b), which computes 3 + 4 $\times$ 2.}
\ans{18}{2}{The dummy int parameter distinguishes the postfix form.}
\ans{19}{2}{The subscript operator [] can be overloaded; .* also cannot be overloaded.}
\ans{20}{4}{The four pillars of object-oriented programming.}
\ans{21}{1}{A pure virtual function is declared as virtual void f() = 0;.}
\ans{22}{2}{Under private inheritance, public and protected members both become private.}
\ans{23}{2}{Virtual inheritance ensures that D contains only one A sub-object.}
\ans{24}{3}{A friend function is not a member; the class grants it access explicitly.}
\ans{25}{1}{Virtual functions provide run-time polymorphism.}
\ans{26}{4}{Static member functions have no this pointer.}
\ans{27}{4}{Constructors cannot be virtual because the object's type is not yet complete during construction.}
\ans{28}{4}{Memory from malloc must be released with free.}
\ans{29}{2}{fstream supports both reading and writing.}
\ans{30}{2}{The compiler may ignore the request.}
\ans{31}{3}{Templates are a compile-time (parametric) polymorphism mechanism.}
\ans{32}{4}{Header files usually hold the extern declarations.}
\ans{33}{2}{C++ has try, throw and catch; finally exists in Java, not in C++.}
\ans{34}{3}{A const member function can be called on both const and non-const objects; inside it, this is a pointer to const.}
\ans{35}{3}{By default terminate calls abort.}
\ans{36}{2}{The framework calls the registered handler when the event is raised.}
\ans{37}{3}{Sq implements the pure virtual function, so it can be instantiated and called through the base pointer.}
\end{answers}

\section{Web Programming (HTML, DHTML, CSS, XML, JavaScript, Java, Servlets and Applets)}

\begin{question}{H}
Match List I with List II.\\
List I (HTML tag): A. <a>  B. <img>  C. <table>  D. <form>\\
List II (Purpose): I. Collects user input to be submitted to a server  II. Creates a hyperlink using the href attribute  III. Embeds an image using the src attribute  IV. Arranges data in rows (<tr>) and cells (<td>)
\optsshort{A-II, B-III, C-I, D-IV}{A-II, B-III, C-IV, D-I}{A-II, B-IV, C-III, D-I}{A-III, B-II, C-IV, D-I}
\end{question}

\begin{question}{H}
In an HTML table, the attribute that makes a cell span two columns is
\optsshort{width="2"}{rowspan="2"}{colspan="2"}{span="2"}
\end{question}

\begin{question}{H}
DHTML (Dynamic HTML) is a combination of
\optslong{XML and XSLT}{HTML and Java applets only}{HTML, CSS, JavaScript and the DOM}{HTTP and FTP}
\end{question}

\begin{question}{H}
Given below are two statements.\\
Statement I: HTML is case-insensitive for tag names, while XML is case-sensitive.\\
Statement II: In XML, every element must be properly closed and nested.
\optslong{Both Statement I and Statement II are true}{Both Statement I and Statement II are false}{Statement I is true but Statement II is false}{Statement I is false but Statement II is true}
\end{question}

\begin{question}{H}
Match List I with List II.\\
List I (XML technology): A. DTD  B. XSD (XML Schema)  C. XSLT  D. XPath\\
List II (Purpose): I. Transforms XML documents into other formats such as HTML  II. Defines document structure with limited data types, in non-XML syntax  III. Selects nodes in an XML document using path expressions  IV. Defines document structure and rich data types, written in XML
\optsshort{A-II, B-I, C-IV, D-III}{A-IV, B-II, C-I, D-III}{A-II, B-IV, C-III, D-I}{A-II, B-IV, C-I, D-III}
\end{question}

\begin{question}{H}
Arrange CSS rules in increasing order of specificity.\\
A. ID selector (\#main)  B. Element selector (p)  C. Inline style attribute  D. Class selector (.note)
\optsshort{B, A, D, C}{D, B, A, C}{B, D, C, A}{B, D, A, C}
\end{question}

\begin{question}{H}
In JavaScript, the values of "5" + 3 and "5" $-$ 3 are respectively
\optsshort{8 and "2"}{"53" and 2}{"53" and "2"}{8 and 2}
\end{question}

\begin{question}{H}
In JavaScript, the results of 0 == "0" and 0 === "0" are respectively
\optsshort{true and false}{true and true}{false and false}{false and true}
\end{question}

\begin{question}{M}
In JavaScript, typeof null evaluates to
\optsshort{"undefined"}{"number"}{"null"}{"object"}
\end{question}

\begin{question}{M}
What is printed by the following JavaScript code?\\
var x = 10; function f() \{ console.log(x); var x = 20; \} f();
\optsshort{20}{ReferenceError}{10}{undefined}
\end{question}

\begin{question}{H}
In Java, with String s1 = "Java", s2 = "Java", s3 = new String("Java");, the values of s1 == s2, s1 == s3 and s1.equals(s3) are
\optsshort{true, true, true}{true, false, false}{false, false, true}{true, false, true}
\end{question}

\begin{question}{H}
In Java, the statements System.out.println("a" + 1 + 2); and System.out.println(1 + 2 + "a"); print
\optsshort{a12 and 12a}{a3 and 3a}{a12 and 3a}{a3 and 12a}
\end{question}

\begin{question}{M}
In Java, System.out.println('a' + 1); prints
\optsshort{Compilation error}{98}{b}{a1}
\end{question}

\begin{question}{M}
In Java, after int x = 5; x += x++ + ++x; the value of x is
\optsshort{16}{12}{17}{18}
\end{question}

\begin{question}{H}
Match List I with List II.\\
List I (Servlet life-cycle method): A. init()  B. service()  C. destroy()\\
List II (When it is called): I. Once, before the servlet is removed from service  II. Once, when the servlet is first loaded  III. For each client request; dispatches to doGet() or doPost()
\optsshort{A-II, B-I, C-III}{A-I, B-III, C-II}{A-III, B-II, C-I}{A-II, B-III, C-I}
\end{question}

\begin{question}{H}
Arrange the life-cycle methods of a Java applet in the order in which they are first called.\\
A. start()  B. init()  C. paint()  D. destroy()  E. stop()
\optsshort{B, C, A, E, D}{B, A, C, E, D}{B, A, C, D, E}{A, B, C, E, D}
\end{question}

\begin{question}{H}
Given below are two statements: one is labelled as Assertion A and the other is labelled as Reason R.\\
Assertion A: Java programs are platform independent.\\
Reason R: Java source code is compiled to native machine code for each processor.
\optslong{Both A and R are true and R is the correct explanation of A}{Both A and R are true but R is NOT the correct explanation of A}{A is true but R is false}{A is false but R is true}
\end{question}

\begin{question}{M}
Which of the following are mechanisms for session tracking in servlets?\\
A. Cookies  B. URL rewriting  C. Hidden form fields  D. HttpSession  E. Applet threads\\
Choose the correct answer from the options given below:
\optsshort{B, C and E only}{A and E only}{A, B, C, D and E}{A, B, C and D only}
\end{question}

\begin{question}{M}
Given below are two statements.\\
Statement I: The HTTP POST method sends form data in the URL.\\
Statement II: The HTTP GET method is preferred for sending large or sensitive data.
\optslong{Both Statement I and Statement II are true}{Both Statement I and Statement II are false}{Statement I is true but Statement II is false}{Statement I is false but Statement II is true}
\end{question}

\begin{question}{M}
Which of the following is true of Java?
\optslong{It supports explicit pointer arithmetic}{It supports multiple inheritance of classes}{It has automatic garbage collection}{It requires manual memory deallocation}
\end{question}

\begin{question}{M}
Match List I with List II.\\
List I (Java keyword): A. final (on a class)  B. static  C. abstract  D. interface\\
List II (Meaning): I. Belongs to the class rather than to instances  II. The class cannot be subclassed  III. Declares a type with abstract methods that a class implements  IV. The class cannot be instantiated and may have abstract methods
\optsshort{A-II, B-I, C-IV, D-III}{A-I, B-II, C-IV, D-III}{A-II, B-IV, C-I, D-III}{A-II, B-I, C-III, D-IV}
\end{question}

\begin{question}{M}
In Java, checked exceptions are those that
\optslong{must be either caught or declared in the throws clause, such as IOException}{occur only at run time}{cannot be caught}{are subclasses of RuntimeException}
\end{question}

\begin{question}{M}
A thread in Java can be created by
\optslong{calling main twice}{using the goto statement}{extending the Thread class or implementing the Runnable interface}{declaring a static block}
\end{question}

\begin{question}{M}
The DOCTYPE declaration <!DOCTYPE html> at the top of a page
\optslong{tells the browser to render the page in HTML5 standards mode}{declares a JavaScript variable}{includes a CSS file}{is an HTML tag that is displayed}
\end{question}

\begin{question}{M}
Given below are two statements.\\
Statement I: An XML document is well-formed if it follows XML syntax rules.\\
Statement II: An XML document is valid only if it is well-formed and conforms to a DTD or schema.
\optslong{Both Statement I and Statement II are true}{Both Statement I and Statement II are false}{Statement I is true but Statement II is false}{Statement I is false but Statement II is true}
\end{question}

\begin{question}{M}
In JavaScript, the event handler attribute that runs a script when a page has finished loading is
\optsshort{onclick}{onchange}{onsubmit}{onload}
\end{question}

\begin{question}{M}
The <meta> tag in HTML is placed in the
\optsshort{<body>}{<table>}{<head>}{<footer>}
\end{question}

\begin{question}{M}
Given below are two statements: one is labelled as Assertion A and the other is labelled as Reason R.\\
Assertion A: XML is suitable for exchanging data between heterogeneous systems.\\
Reason R: XML is a W3C recommendation.
\optslong{Both A and R are true and R is the correct explanation of A}{Both A and R are true but R is NOT the correct explanation of A}{A is true but R is false}{A is false but R is true}
\end{question}

\begin{question}{M}
Which HTML5 element is used for self-contained content such as a blog post?
\optsshort{<frame>}{<article>}{<div>}{<span>}
\end{question}

\begin{question}{L}
CSS stands for
\optsshort{Coded Style Sheets}{Cascading Style Sheets}{Creative Style System}{Computer Style Syntax}
\end{question}

\begin{question}{M}
In JavaScript, the result of 0.1 + 0.2 === 0.3 is
\optsshort{true}{false}{NaN}{undefined}
\end{question}

\begin{question}{M}
JSP (JavaServer Pages) differs from servlets in that JSP
\optslong{replaces HTML entirely}{runs only on the client}{cannot generate dynamic content}{embeds Java code within HTML pages, and is translated into a servlet by the container}
\end{question}

\begin{question}{M}
In Java, the default value of an element of a newly created int array is
\optsshort{undefined}{null}{garbage}{0}
\end{question}

\begin{answers}
\ans{1}{2}{Basic HTML elements.}
\ans{2}{3}{rowspan makes a cell span several rows.}
\ans{3}{3}{Scripts manipulate the DOM and styles to change a page without reloading it.}
\ans{4}{1}{These are among the rules for a well-formed XML document.}
\ans{5}{4}{XML-related W3C standards.}
\ans{6}{4}{!important declarations override normal specificity.}
\ans{7}{2}{+ with a string performs concatenation; $-$ converts its operands to numbers.}
\ans{8}{1}{== converts types before comparing; === compares without type conversion.}
\ans{9}{4}{This is a long-standing quirk of the language; typeof undefined is "undefined".}
\ans{10}{4}{The inner declaration var x is hoisted to the top of f, so x is undefined when logged.}
\ans{11}{4}{Literals are interned and share one object; new creates a separate object; equals compares contents.}
\ans{12}{3}{+ is evaluated left to right; once a String is involved, the rest is concatenation.}
\ans{13}{2}{char is promoted to int, and 'a' is 97.}
\ans{14}{3}{x = 5 + (5 + 7) = 17; Java evaluates operands left to right.}
\ans{15}{4}{The servlet container (such as Tomcat) manages the life cycle.}
\ans{16}{2}{Applets are deprecated in modern Java but remain in the syllabus.}
\ans{17}{3}{Java compiles to bytecode, which runs on any platform that has a JVM.}
\ans{18}{4}{HTTP is stateless, so these mechanisms maintain state across requests.}
\ans{19}{2}{It is the other way round: GET puts data in the URL, and POST (in the request body) suits large or sensitive data.}
\ans{20}{3}{Java supports multiple inheritance of interfaces but not of classes.}
\ans{21}{1}{Basic Java modifiers.}
\ans{22}{1}{Unchecked exceptions, such as NullPointerException, extend RuntimeException.}
\ans{23}{3}{The thread's code goes in the run() method, and start() launches it.}
\ans{24}{1}{Without it, browsers may use quirks mode.}
\ans{25}{1}{Every valid document is well-formed, but not conversely.}
\ans{26}{4}{onclick runs when an element is clicked.}
\ans{27}{3}{It gives metadata such as the character set and description.}
\ans{28}{2}{Both are true, but suitability comes from XML being a self-describing, platform-independent text format.}
\ans{29}{2}{Other semantic elements include <header>, <nav>, <section> and <footer>.}
\ans{30}{2}{CSS separates presentation from content.}
\ans{31}{2}{Binary floating-point cannot represent 0.1 and 0.2 exactly.}
\ans{32}{4}{JSP is convenient for presentation; servlets for control logic.}
\ans{33}{4}{Object reference arrays default to null.}
\end{answers}

\section{Computer Graphics: Video Display Devices, Raster-Scan and Random-Scan Systems, Input Devices}

\begin{question}{H}
The frame buffer size needed for a 1024 $\times$ 768 display with 24 bits per pixel is
\optsshort{3 MB}{1.5 MB}{768 KB}{2.25 MB}
\end{question}

\begin{question}{H}
A raster system with a resolution of 640 $\times$ 480 and 8 bits per pixel needs a frame buffer of
\optsshort{600 KB}{2.4 MB}{300 KB}{307 MB}
\end{question}

\begin{question}{M}
A 1280 $\times$ 1024 display with 24-bit colour needs a frame buffer of
\optsshort{1.25 MB}{5 MB}{3.75 MB}{2.5 MB}
\end{question}

\begin{question}{H}
A raster display of 1024 $\times$ 768 pixels is refreshed 60 times per second. The time available to access each pixel (ignoring retrace) is about
\optsshort{21 ns}{2.1 ns}{16.7 ms}{210 ns}
\end{question}

\begin{question}{H}
With 24 bits per pixel, the number of colours that can be displayed is
\optsshort{about 16.7 million (2\textsuperscript{24})}{24}{65,536}{256}
\end{question}

\begin{question}{H}
A system has an 8-bit frame buffer used as an index into a colour lookup table with 24-bit entries. How many colours can be displayed simultaneously, and from how large a palette?
\optslong{16.7 million out of 256}{256 out of about 16.7 million}{256 out of 256}{8 out of 24}
\end{question}

\begin{question}{H}
Match List I with List II.\\
List I (Display): A. Raster-scan display  B. Random-scan (vector) display  C. Direct-view storage tube  D. Plasma panel\\
List II (Characteristic): I. Draws pictures one line at a time from a display list, producing smooth lines  II. Sweeps the beam row by row and stores pixel intensities in a frame buffer  III. Stores the picture as a charge distribution behind the screen, so no refresh is needed  IV. Flat-panel emissive display using ionised gas between glass plates
\optsshort{A-I, B-II, C-III, D-IV}{A-II, B-I, C-III, D-IV}{A-II, B-I, C-IV, D-III}{A-II, B-III, C-I, D-IV}
\end{question}

\begin{question}{H}
Given below are two statements.\\
Statement I: Random-scan systems produce smooth line drawings because the beam follows the line path directly.\\
Statement II: Random-scan systems are well suited to displaying realistic shaded scenes.
\optslong{Both Statement I and Statement II are true}{Both Statement I and Statement II are false}{Statement I is true but Statement II is false}{Statement I is false but Statement II is true}
\end{question}

\begin{question}{H}
Given below are two statements: one is labelled as Assertion A and the other is labelled as Reason R.\\
Assertion A: Random-scan displays show jagged (aliased) lines.\\
Reason R: In a raster system, lines are approximated by discrete pixel positions on a grid.
\optslong{Both A and R are true and R is the correct explanation of A}{Both A and R are true but R is NOT the correct explanation of A}{A is true but R is false}{A is false but R is true}
\end{question}

\begin{question}{M}
In a CRT, persistence of a phosphor is defined as
\optslong{the number of pixels per inch}{the time to draw one frame}{the time for the emitted light to decay to one-tenth of its original intensity}{the speed of the electron beam}
\end{question}

\begin{question}{M}
Interlaced scanning displays
\optslong{only half the image}{odd-numbered lines in one pass and even-numbered lines in the next}{no image}{all lines in each pass}
\end{question}

\begin{question}{M}
The shadow-mask method for colour CRTs uses
\optslong{a liquid crystal layer}{three electron guns and a grid of red, green and blue phosphor dots}{one electron gun and two phosphor layers}{ionised gas}
\end{question}

\begin{question}{M}
The beam-penetration method of colour CRT display can produce
\optslong{millions of colours}{only four colours: red, orange, yellow and green}{black and white only}{256 colours}
\end{question}

\begin{question}{M}
Given below are two statements.\\
Statement I: An LCD is an emissive display that produces its own light at each pixel.\\
Statement II: In a passive-matrix LCD, a transistor is placed at each pixel position.
\optslong{Both Statement I and Statement II are true}{Both Statement I and Statement II are false}{Statement I is true but Statement II is false}{Statement I is false but Statement II is true}
\end{question}

\begin{question}{M}
Match List I with List II.\\
List I (Input device): A. Light pen  B. Digitizer (graphics tablet)  C. Data glove  D. Trackball\\
List II (Use): I. Grasping and moving virtual objects in three dimensions  II. Pointing directly at positions on the screen by detecting light from the screen  III. Entering coordinates by tracing a drawing on a flat surface  IV. Rotating a ball to move the cursor
\optsshort{A-II, B-III, C-I, D-IV}{A-II, B-III, C-IV, D-I}{A-III, B-II, C-I, D-IV}{A-II, B-I, C-III, D-IV}
\end{question}

\begin{question}{M}
The aspect ratio of a display is
\optslong{the refresh rate}{the ratio of vertical points to horizontal points needed to produce equal-length lines in both directions (or width to height)}{the number of pixels per inch}{the number of colours}
\end{question}

\begin{question}{M}
Resolution of a raster display is
\optslong{the number of colours}{the refresh rate}{the persistence}{the maximum number of points that can be displayed without overlap, often given as pixels across $\times$ down}
\end{question}

\begin{question}{M}
The display processor (graphics controller) in a raster system mainly
\optslong{frees the CPU from graphics tasks such as scan conversion of lines and polygons into the frame buffer}{stores the program}{manages the network}{controls the keyboard}
\end{question}

\begin{question}{M}
Given below are two statements: one is labelled as Assertion A and the other is labelled as Reason R.\\
Assertion A: A random-scan display needs a higher-cost refresh process for complex pictures than for simple ones.\\
Reason R: In a random-scan system, refresh time depends on the number of lines to be drawn.
\optslong{Both A and R are true and R is the correct explanation of A}{Both A and R are true but R is NOT the correct explanation of A}{A is true but R is false}{A is false but R is true}
\end{question}

\begin{question}{M}
A frame buffer that stores one bit per pixel is called a
\optsshort{colour map}{bitmap}{z-buffer}{pixmap}
\end{question}

\begin{question}{L}
Horizontal retrace in a raster CRT is
\optslong{the return of the beam from the bottom to the top of the screen}{the erasing of the screen}{the return of the beam from the right end of a scan line to the left end of the next}{the drawing of a vertical line}
\end{question}

\begin{question}{M}
Given below are two statements.\\
Statement I: A refresh rate below about 24 frames per second causes noticeable flicker.\\
Statement II: Raster-scan systems typically refresh at 60 to 80 frames per second or more.
\optslong{Both Statement I and Statement II are true}{Both Statement I and Statement II are false}{Statement I is true but Statement II is false}{Statement I is false but Statement II is true}
\end{question}

\begin{question}{M}
A raster system has 8 bits per pixel for each of red, green and blue, and a resolution of 512 $\times$ 512. The total frame buffer memory needed is
\optsshort{1 MB}{256 KB}{512 KB}{768 KB}
\end{question}

\begin{question}{M}
In a virtual-reality system, a head-mounted display (HMD) is used to
\optslong{provide a stereoscopic view that changes as the user's head moves}{enter text}{print images}{scan documents}
\end{question}

\begin{answers}
\ans{1}{4}{1024 $\times$ 768 $\times$ 3 bytes = 2,359,296 bytes = 2.25 MB.}
\ans{2}{3}{640 $\times$ 480 = 307,200 bytes = 300 KB.}
\ans{3}{3}{1280 $\times$ 1024 $\times$ 3 bytes = 3,932,160 bytes = 3.75 MB.}
\ans{4}{1}{1 / (1024 $\times$ 768 $\times$ 60) s $\approx$ 21.2 ns.}
\ans{5}{1}{This is called true colour.}
\ans{6}{2}{The lookup table saves frame-buffer memory while allowing a large palette.}
\ans{7}{2}{Common display technologies in Hearn and Baker.}
\ans{8}{3}{Raster systems are needed for realistic shaded and filled images.}
\ans{9}{4}{Random-scan displays draw lines directly and smoothly; jaggedness is a raster-display effect, as R explains.}
\ans{10}{3}{Lower-persistence phosphors need higher refresh rates to avoid flicker.}
\ans{11}{2}{It effectively doubles the refresh rate seen by the eye, reducing flicker at low refresh rates.}
\ans{12}{2}{The beam-penetration method uses two phosphor layers and gives only a few colours.}
\ans{13}{2}{Colour depends on how far the electrons penetrate the red and green layers.}
\ans{14}{2}{An LCD is non-emissive (it modulates light), and per-pixel transistors are the feature of active-matrix LCDs.}
\ans{15}{1}{Graphics input devices.}
\ans{16}{2}{A common screen aspect ratio is 4:3 (or 16:9 for wide screens).}
\ans{17}{4}{For example 1920 $\times$ 1080.}
\ans{18}{1}{Scan conversion converts output primitives into pixel intensities.}
\ans{19}{1}{Too many lines cause flicker because the display cannot be refreshed often enough.}
\ans{20}{2}{With several bits per pixel it is called a pixmap.}
\ans{21}{3}{Vertical retrace returns the beam from the bottom to the top after each frame.}
\ans{22}{1}{Higher refresh rates give a steadier picture.}
\ans{23}{4}{512 $\times$ 512 $\times$ 3 bytes = 786,432 bytes = 768 KB.}
\ans{24}{1}{Tracking devices report the head position to the computer.}
\end{answers}

\section{Output Primitives: Line Drawing (DDA, Bresenham), Mid-Point Circle and Ellipse, Polygon Filling}

\begin{question}{H}
Using Bresenham's line algorithm from (20, 10) to (30, 18), the initial decision parameter p0 = 2$\Delta$y $-$ $\Delta$x is
\optsshort{2}{8}{16}{6}
\end{question}

\begin{question}{H}
Bresenham's line algorithm from (20, 10) to (30, 18) plots which pixel immediately after (22, 12)?
\optsshort{(22, 13)}{(24, 13)}{(23, 13)}{(23, 12)}
\end{question}

\begin{question}{H}
For a line from (20, 10) to (30, 18) drawn with Bresenham's algorithm, the number of pixels plotted, including both endpoints, is
\optsshort{18}{10}{8}{11}
\end{question}

\begin{question}{H}
Using the DDA algorithm with rounding (0.5 rounded up) for the line from (2, 3) to (12, 8), the first four pixels plotted are
\optslong{(2, 3), (3, 3.5), (4, 4), (5, 4.5)}{(2, 3), (3, 4), (4, 4), (5, 5)}{(2, 3), (3, 3), (4, 4), (5, 4)}{(2, 3), (4, 4), (6, 5), (8, 6)}
\end{question}

\begin{question}{H}
Given below are two statements.\\
Statement I: The DDA algorithm uses floating-point arithmetic and rounding at each step.\\
Statement II: Bresenham's line algorithm uses only integer addition, subtraction and multiplication by 2.
\optslong{Both Statement I and Statement II are true}{Both Statement I and Statement II are false}{Statement I is true but Statement II is false}{Statement I is false but Statement II is true}
\end{question}

\begin{question}{H}
In the mid-point circle algorithm for radius r, the initial decision parameter p0 is
\optslong{1 $-$ r (that is, 5/4 $-$ r rounded)}{2r $-$ 1}{r}{r\textsuperscript{2}}
\end{question}

\begin{question}{H}
Using the mid-point circle algorithm for radius 10 centred at the origin, the pixels computed in the first octant are
\optslong{(0, 10), (1, 9), (2, 8), (3, 7), (4, 6), (5, 5)}{(0, 10), (2, 10), (4, 9), (6, 8), (7, 7)}{(0, 10), (1, 10), (2, 10), (3, 10), (4, 9), (5, 9), (6, 8), (7, 7)}{(0, 10), (1, 10), (2, 9), (3, 9), (4, 8), (5, 8), (6, 7), (7, 7)}
\end{question}

\begin{question}{H}
Given below are two statements: one is labelled as Assertion A and the other is labelled as Reason R.\\
Assertion A: The mid-point circle algorithm computes points for only one octant.\\
Reason R: A circle is symmetric only about the x-axis.
\optslong{Both A and R are true and R is the correct explanation of A}{Both A and R are true but R is NOT the correct explanation of A}{A is true but R is false}{A is false but R is true}
\end{question}

\begin{question}{M}
For the mid-point ellipse algorithm with rx = 8 and ry = 6, the initial decision parameter in region 1, p1(0) = ry\textsuperscript{2} $-$ rx\textsuperscript{2} ry + rx\textsuperscript{2}/4, is
\optsshort{332}{$-$300}{$-$332}{$-$348}
\end{question}

\begin{question}{M}
The mid-point ellipse algorithm divides the first quadrant into two regions according to
\optslong{whether the slope of the curve is less than or greater than $-$1 in magnitude}{the sign of x}{the radius}{the colour of the pixels}
\end{question}

\begin{question}{M}
An ellipse has symmetry in
\optsshort{2 halves only}{8 octants}{no symmetry}{4 quadrants}
\end{question}

\begin{question}{H}
Match List I with List II.\\
List I (Fill algorithm or test): A. Scan-line polygon fill  B. Boundary fill  C. Flood fill  D. Odd-even rule\\
List II (Idea): I. Fills an interior region of a uniform colour, replacing it with a new colour  II. Starts at an interior point and fills outward until the boundary colour is reached  III. A point is inside if a ray from it crosses the boundary an odd number of times  IV. Finds intersections of each scan line with polygon edges and fills between pairs
\optsshort{A-II, B-IV, C-I, D-III}{A-IV, B-II, C-III, D-I}{A-IV, B-II, C-I, D-III}{A-IV, B-I, C-II, D-III}
\end{question}

\begin{question}{H}
Given below are two statements.\\
Statement I: A 4-connected boundary fill can leak through diagonal gaps in a boundary that an 8-connected fill would not.\\
Statement II: An 8-connected fill tests the four diagonal neighbours in addition to the four horizontal and vertical neighbours.
\optslong{Both Statement I and Statement II are true}{Both Statement I and Statement II are false}{Statement I is true but Statement II is false}{Statement I is false but Statement II is true}
\end{question}

\begin{question}{M}
In the scan-line polygon fill algorithm, the data structure that holds the edges intersected by the current scan line is the
\optslong{frame buffer}{z-buffer}{active edge table (active edge list)}{edge table (sorted edge table)}
\end{question}

\begin{question}{M}
In the scan-line polygon fill algorithm, a vertex where the two edges lie on opposite sides of the scan line is
\optslong{counted twice}{counted once (often by shortening one of the edges)}{filled with the boundary colour}{ignored}
\end{question}

\begin{question}{M}
The number of pixels plotted by a DDA or Bresenham algorithm for a line from (x1, y1) to (x2, y2) is
\optslong{1}{(x2 $-$ x1)(y2 $-$ y1)}{|x2 $-$ x1| + |y2 $-$ y1|}{max(|x2 $-$ x1|, |y2 $-$ y1|) + 1}
\end{question}

\begin{question}{M}
Using Bresenham's algorithm from (0, 0) to (8, 4), the initial decision parameter is
\optsshort{4}{8}{0}{$-$8}
\end{question}

\begin{question}{M}
Using the mid-point circle algorithm for radius 8, the initial decision parameter is
\optsshort{$-$8}{1}{7}{$-$7}
\end{question}

\begin{question}{M}
Given below are two statements: one is labelled as Assertion A and the other is labelled as Reason R.\\
Assertion A: Recursive flood-fill algorithms can overflow the stack for large regions.\\
Reason R: Each pixel may generate several recursive calls that remain pending.
\optslong{Both A and R are true and R is the correct explanation of A}{Both A and R are true but R is NOT the correct explanation of A}{A is true but R is false}{A is false but R is true}
\end{question}

\begin{question}{M}
Antialiasing techniques such as supersampling and area sampling are used to
\optslong{compress images}{speed up line drawing}{increase resolution in hardware}{reduce the jagged appearance of lines and edges}
\end{question}

\begin{question}{L}
Scan conversion is the process of
\optslong{compressing images}{converting colour to grayscale}{converting an analog image to digital}{converting a geometric primitive, such as a line or circle, into a set of pixels}
\end{question}

\begin{question}{M}
Given below are two statements.\\
Statement I: The non-zero winding number rule and the odd-even rule always give the same result.\\
Statement II: For self-intersecting polygons, the two rules can classify some regions differently.
\optslong{Both Statement I and Statement II are true}{Both Statement I and Statement II are false}{Statement I is true but Statement II is false}{Statement I is false but Statement II is true}
\end{question}

\begin{question}{M}
In the DDA algorithm for a line with slope m where |m| $\leq$ 1, at each step
\optslong{x is incremented by m}{x is incremented by 1 and y by m}{y is incremented by 1 and x by 1/m}{both x and y are incremented by 1}
\end{question}

\begin{answers}
\ans{1}{4}{$\Delta$x = 10 and $\Delta$y = 8, so p0 = 16 $-$ 10 = 6.}
\ans{2}{4}{The decision parameter at that step is $-$2 < 0, so y does not change: (23, 12).}
\ans{3}{4}{For a gentle slope, one pixel is plotted for each x from 20 to 30.}
\ans{4}{2}{The slope is 0.5, so y increases by 0.5 per step: 3.5 rounds to 4, 4.0 stays 4, 4.5 rounds to 5.}
\ans{5}{1}{This is why Bresenham's algorithm is faster and better suited to hardware.}
\ans{6}{1}{For r = 10, p0 = $-$9.}
\ans{7}{3}{This is the standard example in Hearn and Baker.}
\ans{8}{3}{A circle is symmetric in all eight octants, which is why one octant suffices.}
\ans{9}{3}{36 $-$ 384 + 16 = $-$332.}
\ans{10}{1}{In region 1 x is stepped; in region 2, y is stepped.}
\ans{11}{4}{Unlike a circle, an ellipse is not symmetric about the line y = x.}
\ans{12}{3}{The non-zero winding-number rule is another inside-outside test.}
\ans{13}{4}{An 8-connected fill is the one that can leak through diagonal gaps; Statement I reverses the roles.}
\ans{14}{3}{The sorted edge table holds all edges, bucketed by their minimum y.}
\ans{15}{2}{A vertex at a local extremum (both edges on the same side) is counted twice.}
\ans{16}{4}{One pixel per unit step along the major axis, plus the starting pixel.}
\ans{17}{3}{p0 = 2$\Delta$y $-$ $\Delta$x = 8 $-$ 8 = 0.}
\ans{18}{4}{p0 = 1 $-$ r.}
\ans{19}{1}{Scan-line (span) fill methods reduce the stack depth.}
\ans{20}{4}{Pixel intensities are adjusted according to how much of the pixel the line covers.}
\ans{21}{4}{It is also called rasterization.}
\ans{22}{4}{For a five-pointed star, the central pentagon is outside by the odd-even rule but inside by the non-zero rule.}
\ans{23}{2}{For |m| > 1, y is incremented by 1 and x by 1/m.}
\end{answers}

\section{2-D Geometric Transformations (translation, scaling, rotation, reflection, shear, homogeneous coordinates, composite transforms)}

\begin{question}{H}
Rotating the point (3, 4) counter-clockwise by 90$^{\circ}$ about the origin gives
\optsshort{($-$4, 3)}{(4, 3)}{(4, $-$3)}{($-$3, $-$4)}
\end{question}

\begin{question}{H}
Rotating the point (2, 0) counter-clockwise by 30$^{\circ}$ about the origin gives approximately
\optsshort{(1.732, 1)}{(1.732, $-$1)}{(2, 1)}{(1, 1.732)}
\end{question}

\begin{question}{H}
Rotating the point (4, 2) counter-clockwise by 90$^{\circ}$ about the pivot point (2, 2) gives
\optsshort{(0, 2)}{($-$2, 4)}{(2, 4)}{(4, 4)}
\end{question}

\begin{question}{H}
Scaling the point (3, 3) by a factor of 2 in both directions about the fixed point (2, 2) gives
\optsshort{(6, 6)}{(5, 5)}{(8, 8)}{(4, 4)}
\end{question}

\begin{question}{H}
Match List I with List II.\\
List I (Reflection of the point (3, 4)): A. About the x-axis  B. About the y-axis  C. About the line y = x  D. About the origin\\
List II (Result): I. (4, 3)  II. (3, $-$4)  III. ($-$3, $-$4)  IV. ($-$3, 4)
\optsshort{A-II, B-IV, C-I, D-III}{A-IV, B-II, C-I, D-III}{A-II, B-IV, C-III, D-I}{A-II, B-I, C-IV, D-III}
\end{question}

\begin{question}{H}
An x-direction shear with shear factor shx = 2 maps the point (1, 1) to
\optsshort{(2, 2)}{(3, 1)}{(3, 3)}{(1, 3)}
\end{question}

\begin{question}{M}
A y-direction shear with shear factor shy = 3 maps the point (2, 1) to
\optsshort{(5, 1)}{(2, 7)}{(8, 1)}{(2, 4)}
\end{question}

\begin{question}{H}
Homogeneous coordinates are used in 2-D graphics mainly so that
\optslong{curves can be drawn}{translation, as well as rotation and scaling, can be represented as a matrix multiplication, allowing composite transformations to be combined into one matrix}{the image can be stored in less memory}{colours can be represented}
\end{question}

\begin{question}{H}
Given below are two statements.\\
Statement I: Matrix multiplication of 2-D transformations is always commutative.\\
Statement II: Translation cannot be written as a matrix multiplication even in homogeneous coordinates.
\optslong{Both Statement I and Statement II are true}{Both Statement I and Statement II are false}{Statement I is true but Statement II is false}{Statement I is false but Statement II is true}
\end{question}

\begin{question}{H}
Given below are two statements: one is labelled as Assertion A and the other is labelled as Reason R.\\
Assertion A: Two successive translations are multiplicative.\\
Reason R: T(tx1, ty1) $\cdot$ T(tx2, ty2) = T(tx1 + tx2, ty1 + ty2).
\optslong{Both A and R are true and R is the correct explanation of A}{Both A and R are true but R is NOT the correct explanation of A}{A is true but R is false}{A is false but R is true}
\end{question}

\begin{question}{H}
The composite transformation for rotation by $\theta$ about an arbitrary pivot point (xr, yr) is
\optslong{T(xr, yr) $\cdot$ R($\theta$) $\cdot$ T($-$xr, $-$yr)}{T($-$xr, $-$yr) $\cdot$ R($\theta$) $\cdot$ T(xr, yr)}{R($\theta$) only}{R($\theta$) $\cdot$ T(xr, yr)}
\end{question}

\begin{question}{M}
The inverse of a rotation matrix R($\theta$) is
\optsshort{the identity matrix}{undefined}{R($\theta$) itself}{its transpose, R($-\theta$)}
\end{question}

\begin{question}{M}
The 3 $\times$ 3 homogeneous matrix for translation by (tx, ty) has tx and ty in
\optslong{the last column (first two rows), with the remaining part being the identity}{every entry}{the last row}{the diagonal}
\end{question}

\begin{question}{M}
Given below are two statements.\\
Statement I: Uniform scaling preserves the shape (angles) of an object.\\
Statement II: Differential (non-uniform) scaling always preserves the shape of an object.
\optslong{Both Statement I and Statement II are true}{Both Statement I and Statement II are false}{Statement I is true but Statement II is false}{Statement I is false but Statement II is true}
\end{question}

\begin{question}{M}
A scaling transformation with scale factors less than 1
\optslong{enlarges the object and moves it away from the origin}{reduces the object and moves it closer to the coordinate origin}{reflects the object}{rotates the object}
\end{question}

\begin{question}{M}
Reflection about the line y = x is equivalent to
\optslong{a rotation of 90$^{\circ}$}{a translation}{a shear}{a reflection about the x-axis followed by a counter-clockwise rotation of 90$^{\circ}$}
\end{question}

\begin{question}{M}
Rotating the point (1, 0) counter-clockwise by 45$^{\circ}$ about the origin gives approximately
\optsshort{(1, 1)}{(0.5, 0.5)}{(0.707, 0.707)}{(0.707, $-$0.707)}
\end{question}

\begin{question}{M}
Given below are two statements: one is labelled as Assertion A and the other is labelled as Reason R.\\
Assertion A: A rigid-body transformation preserves distances and angles.\\
Reason R: A rigid-body transformation consists only of rotations, translations and reflections.
\optslong{Both A and R are true and R is the correct explanation of A}{Both A and R are true but R is NOT the correct explanation of A}{A is true but R is false}{A is false but R is true}
\end{question}

\begin{question}{M}
The determinant of a 2-D rotation matrix is
\optsshort{1}{0}{cos $\theta$}{$-$1}
\end{question}

\begin{question}{M}
Transforming an object from one Cartesian coordinate system to another, whose origin is at (x0, y0) and whose axes are rotated by $\theta$, is done by
\optsshort{S($\theta$) only}{R($-\theta$) $\cdot$ T($-$x0, $-$y0)}{R($\theta$) $\cdot$ T(x0, y0)}{T(x0, y0) $\cdot$ R($\theta$)}
\end{question}

\begin{question}{M}
The matrix [[1, 0, 0], [0, $-$1, 0], [0, 0, 1]] in homogeneous coordinates represents
\optsshort{reflection about the x-axis}{rotation by 90$^{\circ}$}{reflection about the y-axis}{scaling by $-$1}
\end{question}

\begin{question}{M}
Given below are two statements.\\
Statement I: Shear transformations change the area of an object.\\
Statement II: An affine transformation preserves parallel lines.
\optslong{Both Statement I and Statement II are true}{Both Statement I and Statement II are false}{Statement I is true but Statement II is false}{Statement I is false but Statement II is true}
\end{question}

\begin{question}{L}
In 2-D, a translation moves every point of an object by
\optslong{a reflection}{a scale factor}{a fixed angle}{the same distance in a given direction}
\end{question}

\begin{question}{M}
Scaling an object with sx = 2 and sy = 1 about the origin transforms a unit square with corners (0, 0) and (1, 1) into a rectangle with area
\optsshort{1}{4}{3}{2}
\end{question}

\begin{answers}
\ans{1}{1}{x' = x cos $\theta$ $-$ y sin $\theta$ = $-$4; y' = x sin $\theta$ + y cos $\theta$ = 3.}
\ans{2}{1}{(2 cos 30$^{\circ}$, 2 sin 30$^{\circ}$) = (1.732, 1).}
\ans{3}{3}{Translate by ($-$2, $-$2) to get (2, 0), rotate to (0, 2), and translate back to (2, 4).}
\ans{4}{4}{x' = 2 + 2 $\times$ (3 $-$ 2) = 4, and likewise y' = 4.}
\ans{5}{1}{Reflection about y = $-$x gives ($-$4, $-$3).}
\ans{6}{2}{x' = x + shx $\times$ y = 3; y' = y.}
\ans{7}{2}{y' = y + shy $\times$ x = 1 + 6 = 7; x' = x.}
\ans{8}{2}{A 2-D point (x, y) becomes (x, y, 1), and transformations become 3 $\times$ 3 matrices.}
\ans{9}{2}{Composite transformations are generally not commutative, and homogeneous coordinates exist precisely so that translation becomes a matrix product.}
\ans{10}{4}{R shows that successive translations are additive; successive scalings are multiplicative.}
\ans{11}{1}{Applied right to left: move the pivot to the origin, rotate, then move back.}
\ans{12}{4}{Rotation matrices are orthogonal with determinant 1.}
\ans{13}{1}{Using column vectors, [x', y', 1] = T $\cdot$ [x, y, 1].}
\ans{14}{3}{Non-uniform scaling distorts shapes, for example turning a circle into an ellipse.}
\ans{15}{2}{Scaling about the origin also changes the object's position unless a fixed point is used.}
\ans{16}{4}{Its matrix swaps the x and y coordinates.}
\ans{17}{3}{(cos 45$^{\circ}$, sin 45$^{\circ}$).}
\ans{18}{1}{Scaling and shear are not rigid-body transformations.}
\ans{19}{1}{Rotations preserve area and orientation.}
\ans{20}{2}{First translate the new origin to the old origin, then rotate the axes to align them.}
\ans{21}{1}{It negates the y-coordinate.}
\ans{22}{4}{A shear matrix has determinant 1, so it preserves area while distorting shape.}
\ans{23}{4}{It is a rigid-body transformation.}
\ans{24}{4}{The area is multiplied by |sx $\times$ sy| = 2.}
\end{answers}

\section{2-D Viewing (window-to-viewport mapping, viewing pipeline) and Line and Polygon Clipping}

\begin{question}{H}
A window with corners (0, 0) and (10, 10) is mapped to a viewport with corners (100, 100) and (300, 300). The point (5, 5) maps to
\optsshort{(150, 150)}{(105, 105)}{(250, 250)}{(200, 200)}
\end{question}

\begin{question}{H}
A window with corners (10, 10) and (50, 60) is mapped to the normalised viewport (0, 0) to (1, 1). The point (20, 30) maps to
\optsshort{(0.4, 0.25)}{(0.5, 0.5)}{(0.25, 0.4)}{(0.2, 0.3)}
\end{question}

\begin{question}{H}
In the window-to-viewport transformation, the scaling factor in x is
\optslong{(xvmax $-$ xvmin) / (xwmax $-$ xwmin)}{xvmax / xwmax}{xwmin / xvmin}{(xwmax $-$ xwmin) / (xvmax $-$ xvmin)}
\end{question}

\begin{question}{H}
In the Cohen-Sutherland algorithm with region-code bits ordered Top, Bottom, Right, Left, and a clip window from (0, 0) to (10, 10), the region code of the point ($-$3, 15) is
\optsshort{1001}{0110}{1010}{0101}
\end{question}

\begin{question}{H}
In the Cohen-Sutherland algorithm with bits ordered Top, Bottom, Right, Left and a clip window from (0, 0) to (10, 10), the region code of the point (12, 5) is
\optsshort{0100}{1000}{0001}{0010}
\end{question}

\begin{question}{H}
Given below are two statements.\\
Statement I: In Cohen-Sutherland clipping, a line is trivially accepted if both endpoint codes are 0000.\\
Statement II: In Cohen-Sutherland clipping, a line is trivially rejected if the logical OR of the endpoint codes is non-zero.
\optslong{Both Statement I and Statement II are true}{Both Statement I and Statement II are false}{Statement I is true but Statement II is false}{Statement I is false but Statement II is true}
\end{question}

\begin{question}{H}
A line from ($-$5, 3) to (15, 9) is clipped against the window (0, 0)--(10, 10) by the Liang-Barsky algorithm. The parameter values u1 and u2 of the visible part are
\optsshort{0.25 and 1}{0.25 and 0.75}{0 and 1}{0.5 and 1}
\end{question}

\begin{question}{H}
Match List I with List II.\\
List I (Clipping algorithm): A. Cohen-Sutherland  B. Liang-Barsky  C. Sutherland-Hodgman  D. Weiler-Atherton\\
List II (Characteristic): I. Polygon clipping that also handles concave polygons correctly by following boundaries  II. Line clipping using 4-bit region codes  III. Polygon clipping against each edge of a convex window in turn  IV. Line clipping using the parametric line equation, generally more efficient
\optsshort{A-II, B-IV, C-III, D-I}{A-II, B-IV, C-I, D-III}{A-II, B-III, C-IV, D-I}{A-IV, B-II, C-III, D-I}
\end{question}

\begin{question}{H}
Given below are two statements: one is labelled as Assertion A and the other is labelled as Reason R.\\
Assertion A: The Sutherland-Hodgman algorithm can produce extraneous connecting edges when clipping a concave polygon.\\
Reason R: Sutherland-Hodgman outputs a single vertex list, so separate visible pieces are joined along the window boundary.
\optslong{Both A and R are true and R is the correct explanation of A}{Both A and R are true but R is NOT the correct explanation of A}{A is true but R is false}{A is false but R is true}
\end{question}

\begin{question}{M}
In the Sutherland-Hodgman algorithm, when an edge goes from a vertex inside the clip boundary to a vertex outside it, the algorithm outputs
\optslong{only the intersection point with the boundary}{both vertices}{nothing}{only the inside vertex}
\end{question}

\begin{question}{M}
Arrange the 2-D viewing pipeline stages in order.\\
A. Normalised coordinates  B. Device coordinates  C. Modelling coordinates  D. World coordinates  E. Viewing coordinates
\optsshort{D, C, E, A, B}{C, D, A, E, B}{C, E, D, A, B}{C, D, E, A, B}
\end{question}

\begin{question}{M}
The region of the scene selected for display is called the
\optsshort{frame buffer}{viewport}{window}{clip code}
\end{question}

\begin{question}{M}
Given below are two statements.\\
Statement I: Making the viewport smaller while keeping the window fixed shrinks the displayed picture.\\
Statement II: Making the window larger while keeping the viewport fixed shows a smaller part of the scene in more detail (zooming in).
\optslong{Both Statement I and Statement II are true}{Both Statement I and Statement II are false}{Statement I is true but Statement II is false}{Statement I is false but Statement II is true}
\end{question}

\begin{question}{M}
In Cohen-Sutherland clipping with bits Top, Bottom, Right, Left and the window (0, 0)--(10, 10), a line from ($-$1, $-$1) to (5, $-$2) is
\optslong{clipped against two edges}{trivially accepted}{trivially rejected, because both codes have the bottom bit set}{clipped against the left edge only}
\end{question}

\begin{question}{M}
The Nicholl-Lee-Nicholl line clipping algorithm
\optslong{avoids multiple clipping of an individual line segment by creating more regions around the window}{clips polygons}{clips curves}{uses region codes only}
\end{question}

\begin{question}{M}
In the Liang-Barsky algorithm, a parameter p\_\allowbreak{}k = 0 for some boundary means that the line is
\optslong{parallel to that boundary}{outside the window}{perpendicular to that boundary}{a point}
\end{question}

\begin{question}{M}
Text clipping strategies include
\optslong{only polygon clipping}{all-or-none string clipping, all-or-none character clipping and clipping individual characters}{the z-buffer}{only line clipping}
\end{question}

\begin{question}{L}
Clipping is the process of
\optslong{filling a polygon}{scaling an object}{removing the parts of a picture that lie outside a specified region}{rotating an object}
\end{question}

\begin{question}{M}
Given below are two statements.\\
Statement I: Liang-Barsky clipping generally needs fewer intersection calculations than Cohen-Sutherland clipping.\\
Statement II: Cohen-Sutherland clipping can be extended to three dimensions using 6-bit region codes.
\optslong{Both Statement I and Statement II are true}{Both Statement I and Statement II are false}{Statement I is true but Statement II is false}{Statement I is false but Statement II is true}
\end{question}

\begin{question}{M}
The Weiler-Atherton polygon clipping algorithm differs from Sutherland-Hodgman in that it
\optslong{cannot clip concave polygons}{processes the polygon edges in random order}{clips only lines}{follows the polygon boundary and switches to the window boundary at exit points, producing separate polygons where needed}
\end{question}

\begin{answers}
\ans{1}{4}{sx = sy = 200 / 10 = 20; x' = 100 + (5 $-$ 0) $\times$ 20 = 200.}
\ans{2}{3}{x' = (20 $-$ 10) / 40 = 0.25; y' = (30 $-$ 10) / 50 = 0.4.}
\ans{3}{1}{If sx $\neq$ sy, the object is distorted.}
\ans{4}{1}{The point is above the window (top bit) and to its left (left bit).}
\ans{5}{4}{It lies only to the right of the window.}
\ans{6}{3}{The test for trivial rejection is a non-zero logical AND, meaning both endpoints are outside the same edge.}
\ans{7}{2}{The line enters at x = 0 (u = 5/20) and leaves at x = 10 (u = 15/20); the clipped segment is (0, 4.5) to (10, 7.5).}
\ans{8}{1}{Standard 2-D clipping algorithms.}
\ans{9}{1}{The Weiler-Atherton algorithm avoids this problem.}
\ans{10}{1}{Inside to inside outputs the second vertex; outside to inside outputs the intersection and the second vertex; outside to outside outputs nothing.}
\ans{11}{4}{Clipping is usually performed in normalised coordinates.}
\ans{12}{3}{The viewport is the area of the display device where the window is mapped.}
\ans{13}{3}{A larger window shows more of the scene at a smaller scale (zooming out).}
\ans{14}{3}{The codes are 0101 and 0100; their logical AND, 0100, is non-zero.}
\ans{15}{1}{It needs fewer intersection calculations than Cohen-Sutherland but works only in 2-D.}
\ans{16}{1}{If additionally q\_\allowbreak{}k < 0, the line is completely outside that boundary.}
\ans{17}{2}{They trade accuracy for speed.}
\ans{18}{3}{The region is usually a rectangular clip window.}
\ans{19}{1}{In 3-D, there are six clipping planes, so the codes have six bits.}
\ans{20}{4}{It is suitable for arbitrary (including concave) polygons.}
\end{answers}

\section{3-D Object Representation, Curves and Surfaces, Projections, Illumination and Rendering}

\begin{question}{H}
For the cubic B\'ezier curve with control points P0 (0, 0), P1 (1, 2), P2 (3, 2) and P3 (4, 0), the point at t = 0.5 is
\optsshort{(2, 1.5)}{(2, 1)}{(2, 2)}{(1.5, 2)}
\end{question}

\begin{question}{M}
For the quadratic B\'ezier curve with control points (0, 0), (2, 4) and (4, 0), the point at t = 0.5 is
\optsshort{(2, 2)}{(2, 1)}{(2, 4)}{(4, 2)}
\end{question}

\begin{question}{H}
A B\'ezier curve with n + 1 control points is a polynomial of degree
\optsshort{2n}{3 always}{n}{n + 1}
\end{question}

\begin{question}{H}
Which of the following are properties of B\'ezier curves?\\
A. The curve passes through the first and last control points  B. The curve lies within the convex hull of its control points  C. Moving one control point changes the whole curve  D. The degree is independent of the number of control points\\
Choose the correct answer from the options given below:
\optsshort{A, B, C and D}{A, B and C only}{A, B and D only}{B, C and D only}
\end{question}

\begin{question}{H}
Given below are two statements: one is labelled as Assertion A and the other is labelled as Reason R.\\
Assertion A: B-spline curves provide local control.\\
Reason R: A B-spline curve lies within the convex hull of its control points.
\optslong{Both A and R are true and R is the correct explanation of A}{Both A and R are true but R is NOT the correct explanation of A}{A is true but R is false}{A is false but R is true}
\end{question}

\begin{question}{H}
Match List I with List II.\\
List I (Parallel projection): A. Orthographic  B. Isometric  C. Cavalier  D. Cabinet\\
List II (Characteristic): I. Oblique projection in which lines perpendicular to the view plane are drawn at full length  II. Projectors perpendicular to the view plane, showing front, top or side views  III. Oblique projection in which receding lines are drawn at half length  IV. Axonometric projection in which all three axes are equally foreshortened
\optsshort{A-IV, B-II, C-I, D-III}{A-II, B-IV, C-I, D-III}{A-II, B-IV, C-III, D-I}{A-II, B-I, C-IV, D-III}
\end{question}

\begin{question}{H}
Given below are two statements.\\
Statement I: In a perspective projection, objects farther from the viewer appear smaller.\\
Statement II: A perspective projection preserves the relative proportions of objects regardless of their distance.
\optslong{Both Statement I and Statement II are true}{Both Statement I and Statement II are false}{Statement I is true but Statement II is false}{Statement I is false but Statement II is true}
\end{question}

\begin{question}{H}
A perspective projection with two vanishing points is called
\optsshort{one-point perspective}{two-point perspective}{parallel projection}{three-point perspective}
\end{question}

\begin{question}{H}
Match List I with List II.\\
List I (Shading method): A. Flat (constant) shading  B. Gouraud shading  C. Phong shading\\
List II (Method): I. Interpolates surface normals across the polygon and computes intensity at each pixel  II. Uses a single intensity for each polygon  III. Computes intensities at vertices and interpolates them across the polygon
\optsshort{A-I, B-III, C-II}{A-II, B-III, C-I}{A-II, B-I, C-III}{A-III, B-II, C-I}
\end{question}

\begin{question}{H}
Given below are two statements: one is labelled as Assertion A and the other is labelled as Reason R.\\
Assertion A: Gouraud shading can miss specular highlights that fall inside a polygon.\\
Reason R: Gouraud shading computes intensities only at the vertices and linearly interpolates them.
\optslong{Both A and R are true and R is the correct explanation of A}{Both A and R are true but R is NOT the correct explanation of A}{A is true but R is false}{A is false but R is true}
\end{question}

\begin{question}{H}
In the basic illumination model, diffuse (Lambertian) reflection is proportional to
\optslong{cos $\theta$, where $\theta$ is the angle between the surface normal and the light direction}{cos\textsuperscript{n} $\varphi$, where $\varphi$ is the angle between the reflection and viewing directions}{the area of the polygon}{the distance to the viewer}
\end{question}

\begin{question}{H}
Match List I with List II.\\
List I (Visible-surface method): A. Z-buffer (depth buffer)  B. Painter's (depth-sort) algorithm  C. Back-face detection  D. Warnock's algorithm\\
List II (Idea): I. Recursive subdivision of the screen area until each area is simple  II. Removes polygons whose normal faces away from the viewer  III. Stores the depth of the nearest surface at each pixel  IV. Draws polygons from back to front so nearer ones overwrite farther ones
\optsshort{A-III, B-II, C-IV, D-I}{A-III, B-IV, C-II, D-I}{A-IV, B-III, C-II, D-I}{A-III, B-IV, C-I, D-II}
\end{question}

\begin{question}{H}
Given below are two statements.\\
Statement I: The z-buffer algorithm requires a depth buffer with one entry per pixel in addition to the frame buffer.\\
Statement II: The z-buffer algorithm requires the polygons to be sorted by depth before processing.
\optslong{Both Statement I and Statement II are true}{Both Statement I and Statement II are false}{Statement I is true but Statement II is false}{Statement I is false but Statement II is true}
\end{question}

\begin{question}{M}
In back-face detection with a viewing direction V and polygon normal N (outward), a polygon is a back face if
\optslong{N $\times$ V = 0}{N $\cdot$ V < 0 always}{N $\cdot$ V > 0 (when V points from the viewer into the scene)}{N $\cdot$ V = 0 always}
\end{question}

\begin{question}{M}
Which of the following is a quadric surface?
\optsshort{B\'ezier patch}{Sphere}{Plane}{Polygon mesh}
\end{question}

\begin{question}{M}
Given below are two statements.\\
Statement I: Interpolation splines pass through all their control points.\\
Statement II: Approximation splines, such as B\'ezier and B-spline curves, generally do not pass through all their control points.
\optslong{Both Statement I and Statement II are true}{Both Statement I and Statement II are false}{Statement I is true but Statement II is false}{Statement I is false but Statement II is true}
\end{question}

\begin{question}{M}
C1 (first-order parametric) continuity between two curve sections means that
\optslong{the curves are straight lines}{the curves meet and their first derivatives (tangent vectors) are equal at the joint}{the curves have the same colour}{the curves meet at a common point only}
\end{question}

\begin{question}{M}
Arrange the 3-D viewing pipeline stages in order.\\
A. Projection transformation  B. Modelling transformation  C. Viewing transformation  D. Workstation (viewport) transformation
\optsshort{C, B, A, D}{B, C, A, D}{A, B, C, D}{B, A, C, D}
\end{question}

\begin{question}{M}
A polygon mesh is often stored using
\optslong{a vertex table, an edge table and a polygon (surface-facet) table}{a z-buffer}{a colour table}{only a frame buffer}
\end{question}

\begin{question}{M}
NURBS stands for
\optslong{Normal Uniform Rendering Basis Surfaces}{Non-Uniform Raster Bitmap Shapes}{Non-Uniform Rational B-Splines}{New Uniform Ray-traced B-Surfaces}
\end{question}

\begin{question}{M}
Given below are two statements: one is labelled as Assertion A and the other is labelled as Reason R.\\
Assertion A: Ray tracing can produce reflections, refractions and shadows realistically.\\
Reason R: Ray tracing follows rays from the eye through each pixel into the scene and recursively traces reflected and refracted rays.
\optslong{Both A and R are true and R is the correct explanation of A}{Both A and R are true but R is NOT the correct explanation of A}{A is true but R is false}{A is false but R is true}
\end{question}

\begin{question}{M}
In the 3-D homogeneous representation, a point (x, y, z) is written as
\optsshort{(x, y, 1)}{(x, y, z)}{(1, x, y, z)}{(x, y, z, 1)}
\end{question}

\begin{question}{M}
The 3-D rotation of a point about the z-axis by $\theta$
\optslong{is a translation}{changes x and y as in 2-D rotation and leaves z unchanged}{changes only the z-coordinate}{changes only x}
\end{question}

\begin{question}{M}
Given below are two statements.\\
Statement I: The A-buffer extends the z-buffer to handle transparency and antialiasing.\\
Statement II: The painter's algorithm always handles cyclically overlapping polygons without splitting them.
\optslong{Both Statement I and Statement II are true}{Both Statement I and Statement II are false}{Statement I is true but Statement II is false}{Statement I is false but Statement II is true}
\end{question}

\begin{question}{L}
The phenomenon of light reflected from many surfaces in a scene, modelled crudely as a constant, is called
\optsshort{ambient light}{refraction}{diffuse reflection}{specular reflection}
\end{question}

\begin{question}{M}
In 3-D clipping against a normalised view volume, the Cohen-Sutherland region codes use
\optsshort{6 bits}{3 bits}{8 bits}{4 bits}
\end{question}

\begin{question}{M}
A B\'ezier surface is defined by
\optslong{a single control point}{a z-buffer}{a 1-D array of control points}{a 2-D grid (mesh) of control points, using two parameters u and v}
\end{question}

\begin{answers}
\ans{1}{1}{B(0.5) = 0.125 P0 + 0.375 P1 + 0.375 P2 + 0.125 P3 = (2, 1.5).}
\ans{2}{1}{B(0.5) = 0.25 P0 + 0.5 P1 + 0.25 P2 = (2, 2).}
\ans{3}{3}{Four control points give a cubic B\'ezier curve.}
\ans{4}{2}{D is a property of B-splines, not of B\'ezier curves.}
\ans{5}{2}{Both are true, but local control comes from each basis function being non-zero over only a limited parameter range.}
\ans{6}{2}{Dimetric and trimetric are other axonometric projections.}
\ans{7}{3}{Parallel projection preserves relative proportions; perspective gives realism but not true measurements.}
\ans{8}{2}{The number of principal vanishing points equals the number of principal axes cut by the projection plane.}
\ans{9}{2}{Phong shading renders specular highlights more accurately but costs more.}
\ans{10}{1}{Phong shading interpolates normals instead, so highlights inside polygons appear.}
\ans{11}{1}{Specular reflection in the Phong model is proportional to cos\textsuperscript{n} $\varphi$.}
\ans{12}{2}{Z-buffer is an image-space method; depth sorting combines object and image space.}
\ans{13}{3}{The z-buffer processes polygons in any order; that independence from sorting is its main advantage.}
\ans{14}{3}{Conventions differ; with V pointing from the surface towards the viewer, back faces have N $\cdot$ V < 0.}
\ans{15}{2}{Quadric surfaces (sphere, ellipsoid, torus, paraboloid, hyperboloid) are described by second-degree equations.}
\ans{16}{1}{Cubic B\'ezier curves pass only through the first and last control points.}
\ans{17}{2}{C0 continuity means only that the curves meet; C2 also matches second derivatives.}
\ans{18}{2}{Modelling to world coordinates, world to viewing, projection, then mapping to the device.}
\ans{19}{1}{Geometric tables are supplemented by attribute tables for colour and transparency.}
\ans{20}{3}{NURBS can represent conic sections such as circles exactly.}
\ans{21}{1}{It is computationally expensive.}
\ans{22}{4}{3-D transformations are then 4 $\times$ 4 matrices.}
\ans{23}{2}{Rotations about the x- and y-axes are obtained by cyclic permutation of the coordinates.}
\ans{24}{3}{Cyclic overlap requires splitting polygons in the depth-sort method.}
\ans{25}{1}{Ambient light gives surfaces a base level of illumination.}
\ans{26}{1}{There is one bit for each of the six clipping planes.}
\ans{27}{4}{It is the Cartesian product of B\'ezier blending functions in u and v.}
\end{answers}
